%Planning to understand/simplify/improve. If you are interested, please contact me (Evangelos Kipouridis), I would be more than happy to do that collaboratively.
%Licensed under CC BY 4.0
\documentclass[11pt]{article}
\usepackage[T1]{fontenc}
\usepackage[margin=1in]{geometry}
\usepackage{amsmath,amssymb,amsthm,mathtools,booktabs,enumitem,microtype}
\microtypesetup{expansion=false}
\usepackage[hidelinks]{hyperref}
\newtheorem{theorem}{Theorem}[section]
\newtheorem{lemma}[theorem]{Lemma}
\newtheorem{proposition}[theorem]{Proposition}
\newtheorem{corollary}[theorem]{Corollary}
\theoremstyle{definition}
\newtheorem{definition}[theorem]{Definition}
\newtheorem{remark}[theorem]{Remark}
\newcommand{\E}{\mathbb E}
\newcommand{\Prb}{\mathbb P}
\newcommand{\ind}{\mathbf 1}
\newcommand{\OPT}{\operatorname{OPT}}
\newcommand{\cost}{\operatorname{cost}}
\newcommand{\wt}{\widetilde}
\newcommand{\Part}{\mathcal P}
\newcommand{\join}{\chi}
\newcommand{\cut}{\partial}
\newcommand{\pos}[1]{(#1)_+}
\DeclareMathOperator{\diag}{diag}
\DeclareMathOperator{\MSF}{MSF}
\title{Algorithms for Ultrametric Fitting\\
\large Offline Algorithms and One-Pass Algorithms}
\author{}
\date{29 October 2026}
\begin{document}
\maketitle

\begin{abstract}
An ultrametric is a symmetric nonnegative function with zero diagonal whose value on any pair
is at most the maximum of its values on the other two pairs in any triple.
Given nonnegative distances on all pairs of $n$ vertices, we fit an
ultrametric by minimizing either the sum of absolute errors or the square
root of the sum of squared errors. We give four algorithm families:
offline and one-pass algorithms for each objective. For absolute error,
the offline approximation factors are $O(\log^2 n)$ in
$\widetilde O(n^{10})$ time and $O(\log^5 n)$ in
$\widetilde O(n^2)$ time. The one-pass approximation factors are
$O(\log^3 n)$ with polynomial postprocessing and $O(\log^5 n)$ with
$\widetilde O(n^2)$ postprocessing. For the square root of squared error,
both approximation factors are $O(\log^4 n)$; the offline time is
$\widetilde O(n^2)$ and the one-pass postprocessing time is
$\widetilde O(n^3)$. Each one-pass algorithm stores
$\widetilde O(n)$ words of $O(\log n)$ bits. An offline algorithm may
access any input pair; a one-pass algorithm reads every input pair once,
in an arbitrary fixed order. The notation $\widetilde O$ suppresses powers
of $\log(n+1)$. The computational
assumptions and success probabilities are specified in the text.
\end{abstract}

\section{Definitions and main theorem}
\label{sec:definitions}

For a finite set $X$, let $\binom X2$ denote the set of its two-element
subsets. All sums indexed by $xy\in\binom X2$ count an unordered pair once.
Within a specified ambient set $X$, a sum indexed only by $xy$ means a
sum over $\binom X2$. The ambient set is $V$ unless a state specifies $X$. Every finite indexing set has a fixed total order; newly created atoms receive consecutive identifiers. Ties in minimization are resolved by the resulting lexicographic order unless another rule is specified.
Let $n\ge2$, let $V=\{1,\ldots,n\}$, and let
$d:\binom V2\to\mathbb R_{\ge0}$. Extend $d$ symmetrically to $V^2$ and
put $d_{xx}=0$. Set $M=\max_{xy\in\binom V2}d_{xy}$.
For a statement $A$, $\ind[A]$ is one if $A$ holds and zero otherwise.
For a real number $z$, $z_+=\max\{z,0\}$.
The notation $\log$ denotes the natural logarithm; $\log_2$ denotes the
base-two logarithm. We write $\Prb$ for probability and $\E$ for
expectation. A subscript specifies the random variable being averaged
when it is not clear from the expression.

An \emph{ultrametric} on $X$ is a symmetric map $u:X^2\to\mathbb R_{\ge0}$
such that $u_{xx}=0$ and
\[
 u_{xy}\le\max\{u_{xz},u_{zy}\}\qquad(x,y,z\in X).
\]
Zero distance between distinct elements is permitted. Define
\begin{align*}
 L_1(d,u)&=\sum_{xy\in\binom V2}|d_{xy}-u_{xy}|,\\
 L_2(d,u)&=\sum_{xy\in\binom V2}(d_{xy}-u_{xy})^2,\\
 \OPT_p&=\min_{u\text{ ultrametric}}L_p(d,u)\qquad(p=1,2).
\end{align*}
An $r$-approximation for $L_p$, where $r\ge1$, is an output ultrametric
$u$ satisfying $L_p(d,u)\le r\OPT_p$. For symmetric matrices $A,B$ on
the same ordered vertex set, define
$\|A-B\|_2=(\sum_{x<y}(A_{xy}-B_{xy})^2)^{1/2}$.
The $\ell_2$ norm error is $L_2(d,u)^{1/2}$. An $r$-approximation for
$L_2$ is therefore an $r^{1/2}$-approximation for the norm error.
Replacing every entry of $u$ by $\min\{u_{xy},M\}$ preserves the
ultrametric inequalities and cannot increase either objective. The feasible
set in $[0,M]^{\binom n2}$ is closed and bounded, so each minimum exists.
The height-$M$ ultrametric has off-diagonal entries $M$ and diagonal entries zero. If $M=0$, the zero ultrametric is an exact answer. Henceforth $M>0$.

A \emph{partition} of $X$ is a set of nonempty disjoint subsets whose union
is $X$. Its elements are its \emph{blocks}. For a partition $P$, $P(x)$ is
the block containing $x$ and $\join_P(x,y)=\ind[P(x)=P(y)]$.
$P$ \emph{refines} $Q$ when every block of $P$ is contained in a block of
$Q$. A sequence is \emph{nested} when each partition refines every later
partition. For $t\ge0$, the threshold partition $P_u(t)$ of an ultrametric $u$ is defined
by $\join_{P_u(t)}(x,y)=\ind[u_{xy}\le t]$.

A \emph{hierarchy representation} of $u$ is a rooted tree whose leaves are
$X$, with each internal vertex assigned a nonnegative height, heights
nondecreasing towards the root, and $u_{xy}$ equal to the height of the
lowest common ancestor of $x,y$. The lowest common ancestor is the common
ancestor farthest from the root. Suppress internal vertices with one child.
The representation then has at most $2|X|-1$ vertices. Every ultrametric has
such a representation: adjoin all singleton sets to the family of its
threshold blocks. Two sets in this family are disjoint or one contains
the other. Connect each nonroot set to its least strict superset, and
assign each nonsingleton set the first threshold at which it is a block.
Assign singleton leaves height zero. This also represents distinct
vertices at distance zero, whose common ancestor has height zero. Equal-height internal
vertices may be inserted without changing $u$. All algorithms output a
hierarchy representation, rather than its quadratic table of pair distances.

An \emph{offline algorithm} has random access to the input table.
A \emph{one-pass algorithm} reads one record $(x,y,d_{xy})$ for each
$xy\in\binom V2$, in an arbitrary fixed order independent of its randomness,
and cannot read a record again. Its \emph{postprocessing} is its computation
after the last record. Space includes postprocessing space and the output.
The vertex set is known before the pass. An input prefix is the first
$j$ records for some $j\ge0$. An update is the computation performed
when one record arrives; total update time sums these times over the pass.

Time bounds count word operations. For these bounds, input distances are
nonnegative integers at most $n^b$, for a fixed constant $b$; a common scaling
of all distances is allowed. A word has $O(\log n)$ bits. Numbers occupying
a polylogarithmic number of words are processed in a polylogarithmic number
of word operations. The notation $\wt O(f(n))$ means
$O(f(n)\log^c(n+1))$ for some constant $c$. A probability bound holds
\emph{with high probability} if, for every fixed $a>0$, its failure
probability can be made at most $n^{-a}$ by changing constants in the
algorithm. Unless expressly stated otherwise, all random choices described
as fresh are independent of the previous choices, conditional on the current
input to that step. The word \emph{combinatorial} here means that the
algorithm uses arithmetic, comparison, sampling, graph searches, and finite
dynamic programming, without a linear-programming, semidefinite-programming,
or numerical linear-system subroutine. The faster cut subroutines below are
not asserted to be combinatorial.
We call a quantity polynomially bounded if it is at most $n^c$ for a
fixed constant $c$, and polylogarithmically bounded if it is at most
$\log^c(n+1)$ for a fixed constant $c$. ``Near-linear time'' on a graph
with $N$ vertices and $m$ edges means $\widetilde O(N+m)$ time, with
logarithmic factors measured against the original input size $n$.

\begin{theorem}[The four algorithm families]\label{thm:four}
For complete, unweighted input pairs, the following randomized algorithms
return the indicated approximation with high probability.
The four families are offline $L_1$, one-pass $L_1$, offline $L_2$, and
one-pass $L_2$. The first two families each have two implementations.
An unweighted input pair has coefficient one in each objective.
\begin{center}
\small
\begin{tabular}{@{}llll@{}}\toprule
Objective and model & Approximation & Time & Space\\\midrule
$L_1$, combinatorial offline & $O(\log^2 n)$ & $\wt O(n^{10})$ & polynomial\\
$L_1$, faster offline & $O(\log^5 n)$ & $\wt O(n^2)$ & $\wt O(n^2)$\\
$L_1$, combinatorial one pass & $O(\log^3 n)$ & polynomial & $\wt O(n)$\\
$L_1$, faster one pass & $O(\log^5 n)$ & $\wt O(n^2)$ & $\wt O(n)$\\
$L_2$, offline & $O(\log^8 n)$ & $\wt O(n^2)$ & $\wt O(n^2)$\\
$L_2$, one pass & $O(\log^8 n)$ & $\wt O(n^3)$ & $\wt O(n)$\\\bottomrule
\end{tabular}
\end{center}
For the faster $L_1$ one-pass algorithm, both total update time and
postprocessing time are $\wt O(n^2)$. For the $L_2$ one-pass algorithm,
total update time is $\wt O(n^2)$ and postprocessing time is $\wt O(n^3)$.
The last two approximation factors are $O(\log^4 n)$ for the norm error.
The faster bounds use the capacitated cut algorithms stated in
Section~\ref{sec:primitives}. The proofs of the reductions to those
algorithms are given in full below.
\end{theorem}

\section{Graph algorithms used as external results}
\label{sec:primitives}

A graph in this paper is finite and undirected. Parallel edges are permitted.
An edge capacity is nonnegative. Self-loops may be deleted. For
$J=(Y,E,c)$ and $S\subseteq Y$, put
\[
 \cut_J S=\{xy\in E:|\{x,y\}\cap S|=1\},\qquad
 c_J(\cut S)=\sum_{e\in\cut_J S}c_e.
\]
For $q:Y\to\mathbb Z_{\ge0}$, write $q(S)=\sum_{x\in S}q_x$.
A cut is \emph{pinned} at $s,t$ if $s\in S$ and $t\notin S$. For $q(Y)=N$, call $S$ \emph{narrow-balanced} if $N/3\le q(S)\le2N/3$ and \emph{wide-balanced} if $N/8\le q(S)\le7N/8$.
Contracting a vertex set means replacing it by one vertex, adding capacities
of parallel edges and deleting self-loops. Sum $q$ over the contracted set.
A vertex measure is a function assigning a nonnegative real number to each
vertex. For any vertex measure $q$, the notation $q(S)$ has the same
sum meaning as above. The induced subgraph $J[Z]$, for $Z\subseteq Y$,
retains the vertices of $Z$ and all edges with both endpoints in $Z$.
The minimum cut value of a graph with at least two vertices is the least
capacity of a cut with both sides nonempty. A cut sparsifier with error
$\eta$ is a graph on the same vertices whose capacity on every cut is
between $1-\eta$ and $1+\eta$ times the original capacity. The word
compression will mean replacing a graph by such a graph. A terminal in a
tree is one of the designated leaves corresponding to vertices of the
original graph; all other tree vertices are auxiliary vertices. A clique
on a vertex set contains one edge for each pair of distinct vertices.

\begin{theorem}[External cut algorithms]\label{thm:external}
We use the following results, with inverse-polynomial failure probability.
\begin{enumerate}[label=(\alph*)]
\item \label{ext:tree}
For a graph $J$ on $N$ vertices with $m$ edges, a combinatorial algorithm
constructs a distribution of capacitated trees with $O(N)$ vertices and
leaves $Y$. For a tree $T$, define
\[
 c_T(S)=\min\{c_T(\cut_T Z):Z\subseteq V(T),\ Z\cap Y=S\}.
\]
Every tree in the distribution satisfies $c_J(\cut S)\le c_T(S)$, and
$\E_Tc_T(S)\le\beta(N)c_J(\cut S)$ for each $S\subseteq Y$, where
$\beta(N)=O(\log(N+1))$. A conservative construction time is
$\wt O(N(N+m)^2)$. This is the cut consequence of R\"acke's
Theorems 2 and 4 and Lemmas 7 and 8~\cite{Racke08}, using
the metric-tree construction of~\cite{FRT04}.
\item \label{ext:ed}
For a capacitated graph and a vertex measure $q$, a near-linear-time
algorithm, given $\phi>0$, partitions $Y$ into sets $Y_i$ such that
every $J[Y_i]$ satisfies
\[
 c_{J[Y_i]}(\cut S)\ge a_0\phi\min\{q(S),q(Y_i\setminus S)\}
 \quad(S\subseteq Y_i),
\]
and the capacity of edges between sets is at most
$C_0\phi q(Y)\log^2(N+1)$. Here $a_0,C_0>0$ are absolute constants.
Zero values of $q$ are allowed, and its positive values lie between inverse
polynomial and polynomial bounds. This is Theorem 3.2
of Agassy--Dorfman--Kaplan~\cite{ADK25}, with its weighted-graph conventions.
\item \label{ext:st}
A pinned cut of capacity at most twice the minimum pinned capacity is
computable in $\wt O(N+m)$ time. We use Theorem 1.3
of Li--Nanongkai--Panigrahi--Saranurak~\cite{LNPS23} with constant accuracy.
\item \label{ext:bk}
For $0<\eta<1/2$, a capacitated graph can be replaced, in
$\wt O(N+m)$ time for inverse-polylogarithmic $\eta$, by a graph on the same
vertices with $\wt O(N)$ edges such that all cut capacities change by factors
in $[1-\eta,1+\eta]$. We use Bencz\'ur--Karger~\cite{BK15},
Theorems 6.4--6.5. Its sampling theorem also implies: in an unweighted
multigraph, independently retaining edge $e$ with probability
$p_e\ge\min\{1,\rho/k_e\}$ and assigning retained capacity $1/p_e$
preserves all cuts within $1\pm\eta$, where
$\rho=C\eta^{-2}\log(N/\zeta)$, failure probability is at most $\zeta$,
and $k_e$ is the largest minimum cut of a vertex-induced subgraph containing
the endpoints of $e$. Only the cut guarantee, not a size guarantee, is
claimed for oversampling. This is the compression theorem and its
oversampling consequence, Theorems 5.2 and 6.2 of~\cite{BK15}.
\end{enumerate}
\end{theorem}

The tree construction time follows from $O((N+m)\log(N+1))$ calls to the
metric-tree procedure in Lemmas 7--8 of~\cite{Racke08}. With an explicit
distance table, one call runs shortest-path searches from every source
and constructs a metric tree in $\wt O(N(N+m))$ time. A shortest path
minimizes the sum of its edge lengths, and its vertex-pair minima define
the distance table. A metric tree here is a tree with nonnegative edge
lengths used to approximate that table; it is an intermediate object of
the external construction.

For completeness, the cut formulation in part~\ref{ext:tree} follows
from the flow formulation in~\cite{Racke08} as follows. A path flow from
$S$ to $Y\setminus S$ assigns nonnegative amounts to paths with one
endpoint on each side. Its value is the sum of these amounts, and an
edge's load is the sum of amounts on paths using it. Congestion is the
maximum load divided by capacity; an edge of capacity zero must have
zero load. The maximum-flow minimum-cut theorem identifies the largest
value of a congestion-one path flow with the minimum capacity of a
separating cut. The cited tree maps transfer every graph flow to each
tree without increasing congestion. Thus $c_J(\cut S)\le c_T(S)$.
In the other direction, choose a maximum $S$-to-$Y\setminus S$ flow
in each tree. The cited averaged maps send these flows to $J$ with
congestion at most $\beta(N)$ and total value $\E_Tc_T(S)$.
Every resulting path crosses $\cut_J S$, so its total value is at most
$\beta(N)c_J(\cut S)$. This proves the expectation inequality.
Disconnected positive-capacity components are handled separately and their
trees joined by zero-capacity edges.

\begin{lemma}[Constrained tree cuts]\label{lem:tree-dp}
Given a tree of part~\ref{ext:tree}, pins, integer terminal masses with
total at most $n$, an exact prescribed selected mass, and a prescribed
interval of selected-terminal counts, the minimum tree cut is computable
in $O(n^2N^3)$ time. If there is no mass constraint, the computation and
recovery of one minimizing cut use $O(N)$ working words and polynomial time.
Consequently, with high probability, an $O(\log n)$-approximate pinned
cut with counts in $[N/3,2N/3]$ can be found combinatorially in polynomial time and
$\wt O(N+m)$ space, without an exact-mass constraint.
\end{lemma}
\begin{proof}
Root the tree and, if necessary, replace a vertex of high degree by a binary
tree of copies. Joining copies with edges larger than the sum of all finite
capacities forces identical labels. For a vertex $v$, let
$A_v(\sigma,j,h)$ be the minimum cost inside its subtree when $v$ has
label $\sigma\in\{0,1\}$, the number of terminals with label one is $j$,
and their mass is $h$. A terminal contributes its count and mass precisely
when its label is one; an internal vertex contributes zero. A pin prohibits
the other label. For a child $z$, combine a parent entry and a child entry,
add their count and mass indices, and add the edge capacity exactly when
their labels differ. Minimize over entries with the same resulting indices.
More explicitly, let $a_v$ equal one at a terminal and zero otherwise,
and let $m_v$ be its terminal mass, or zero at an auxiliary vertex.
If $z_1,z_2$ are the two children and their edge capacities are
$c_1,c_2$, the recurrence is
\[
 A_v(\sigma,j,h)=
 \min_{\substack{j_1+j_2=j-a_v\sigma\\h_1+h_2=h-m_v\sigma}}
 \ \sum_{i=1}^2\min_{\sigma_i\in\{0,1\}}
 \bigl(A_{z_i}(\sigma_i,j_i,h_i)+c_i\ind[\sigma_i\ne\sigma]\bigr).
\]
An absent child contributes zero with count and mass zero. A leaf has
value zero exactly at its permitted label and corresponding count and
mass, and infinity otherwise. There are $O(N)$ combinations, and direct
enumeration of two count indices and two mass indices costs
$O(N^2n^2)$ at each combination. This proves the running time.
At the root minimize over permitted indices and both labels. An infeasible
entry has value $+\infty$.

Without the mass index, live arrays have total length $O(N)$ if children are
processed one at a time and arrays are allocated only for processed
terminals. The processed terminal sets underlying simultaneously retained
arrays are disjoint; zero-terminal arrays contribute at most the number
of tree vertices. Discard child arrays after combination. Recover a cut by
fixing vertex labels successively, recomputing the optimum with each next
label forced to zero; keep it if the optimum is unchanged, and otherwise
force one. There are $O(N)$ recomputations and the space bound is unchanged.

For a fixed optimal feasible set $S^*$, a tree minimizer $S_T$ satisfies
$c_J(\cut S_T)\le c_T(S_T)\le c_T(S^*)$. Part~\ref{ext:tree} and Markov's
inequality give a $2\beta(N)$ approximation with probability at least
one half. Repeat independently $O(\log n)$ times and retain the least
original graph cost. If the optimum is zero, its nonnegative tree costs
have expectation zero and vanish with probability one.
The construction in part~\ref{ext:tree} admits $\wt O(N+m)$ workspace:
retain the current multiplicative edge weights and a weighted reservoir
sample of the generated trees. Specifically, if successive trees have
positive weights $p_1,p_2,\ldots$, replace the currently stored tree by tree $i$
with probability $p_i/(p_1+\cdots+p_i)$. Induction shows that after step
$i$ each earlier tree $j$ is stored with probability
$p_j/(p_1+\cdots+p_i)$. Only one tree and the cumulative weight are
stored. Shortest-path arrays for the metric-tree
procedure can be recomputed one source at a time, and each embedded tree-edge
path can be processed and discarded. The metric-tree partitions require
only the current vertex memberships and a permutation; distances are
recomputed on demand. This increases time only polynomially. No full
distance table, collection of embeddings, or collection of trees is stored.
\end{proof}

\paragraph{Coefficient encoding.}
The graph algorithms are used with finite precision in their stated
word-operation models. The reductions do not require exact arithmetic on
unbounded denominators. One may round every positive generated capacity
upwards to a binary number within a factor $1+\nu$, retaining zero
exactly, where $\nu=n^{-c}$ and $c$ is a sufficiently large fixed constant.
For any cut, the same relative bound holds after summing capacities.
Thus an $\alpha_0$-approximate cut for the rounded graph is an
$(1+\nu)\alpha_0$-approximate cut for the unrounded graph.
In cumulative compression, only $O(\log n)$ such factors occur on one
entry's history. Allocate these errors within the constant slack in
Lemma~\ref{lem:prefix-core} and the squared-error graph bounds below.
Use a power-of-two upper bound for each approximation parameter and
choose sign-projection dimension to be a power of two. Penalty scans
then also have binary coefficients after this rounding.

All distances, masses, heights, and uncompressed certificate coefficients
have polynomial magnitude. Product-clique sampling creates positive
capacities at least an inverse polynomial: its sample count and every
probability denominator are polynomial, and its unscaled coefficients
are positive integers. Compression resamples a retained edge with
capacity divided by a probability at most one, so it does not decrease
this lower bound. On the cut-preservation event, summing singleton cuts
bounds total compressed capacity by a constant times original total
capacity; in particular its magnitude stays polynomial. A larger
capacity is detected and flags failure. The prefix sketch capacities
are powers of two with logarithmic exponents. In squared-error states,
all queried entries and heights are integers, and projected squared norms
have denominator equal to the sign dimension before their fixed rescaling.
These facts give a fixed polynomial upper and inverse-polynomial lower
bound for every positive reduction coefficient. Binary rounding therefore
uses $O(\log n)$ exponent and mantissa bits. All positivity tests and
zero-cost cases remain exact. Factors $1+\nu$ are absorbed in the fixed
constants of the cut guarantees, not in the approximation exponents.

\begin{lemma}[Fast pinned balanced cuts]\label{lem:fast-cut}
Let terminal count be one for each of $N$ atoms and zero for one additional
vertex. Pins may force one atom inside and any specified atoms outside.
There is a randomized $\wt O(N+m)$-time algorithm that, with high
probability, returns a pinned cut with between $N/8$ and $7N/8$ atoms,
of capacity at most $\alpha(n)=O(\log^2 n)$ times
every pinned cut with between $N/3$ and $2N/3$ atoms, whenever one exists.
\end{lemma}
\begin{proof}
First consider ordinary count weights $q$ with $q(Y)=N$.
If a $[N/3,2N/3]$-balanced optimum has positive cost $h$, scan geometric
values of $\phi$ including one in $[24h/(a_0N),48h/(a_0N)]$.
Part~\ref{ext:ed} cannot return a component of count greater than $3N/4$:
both its intersections with the comparison cut have count greater than
$N/12$, contradicting expansion and the cut's cost $h$.
Choose a component of count in $[N/4,3N/4]$ if there is one; otherwise
accumulate components of count less than $N/4$ until the count first reaches
$N/4$. The count is less than $N/2$, and the cut costs $O(h\log^2 n)$.
When $h=0$, do this test on the positive-capacity connected components
before the scan. A positive cut value is between the smallest positive
capacity and the sum of capacities, so the parameter scan has
polylogarithmic length under the encoding convention. Explicitly check
balance and retain the least valid result.

Contract the forced-out atoms into the additional vertex $t$, retaining
their count. Write $s$ for the inside pin. Find the ordinary balanced cut
$D$ just described and a twice-optimal pinned cut $A$ by part~\ref{ext:st}.
If $D$ separates $s,t$, orient it towards $s$ and return it. Otherwise orient
$D$ so it contains neither. If $q(A)\in[N/8,7N/8]$, return $A$.
If $q(A)<N/8$, return $A\cup D$; if $q(Y\setminus A)<N/8$, return
$A\setminus D$. Each result has the required pins and count.
For any sets $A,D$, the cut of either $A\cup D$ or $A\setminus D$ has
capacity at most $c_J(\cut A)+c_J(\cut D)$, since every crossing edge
crosses at least one of the two original cuts. Both summands are bounded
against each permitted pinned comparison cut. Repeat the ordinary cut
construction and the pinned cut construction independently $O(\log n)$
times. Check the count constraints for ordinary cuts and the pins for
pinned cuts; in each construction retain the least capacity among valid
outputs. Each required comparison has constant success probability in
one trial, so repetition makes its failure inverse polynomial in $n$.
This proves the claim.
\end{proof}

We use one matrix concentration inequality in the next proof. A real
symmetric matrix $A$ is positive semidefinite when $z^{\mathsf T}Az\ge0$
for every real vector $z$. Write $A\preceq B$ when $B-A$ is positive
semidefinite, and write $I_\ell$ for the identity matrix in dimension
$\ell$. The following form of the matrix Chernoff inequality follows
from Corollary 5.2 of~\cite{Tropp12}: if independent random symmetric
matrices $Y_j$ satisfy $0\preceq Y_j\preceq R_0I_\ell$ and
$\sum_j\E Y_j=I_\ell$, then for $0<\eta<1$,
\[
 \Prb\!\left[(1-\eta)I_\ell\preceq\sum_jY_j
                    \preceq(1+\eta)I_\ell\right]
 \ge1-2\ell\exp\!\left(-\frac{\eta^2}{3R_0}\right).
\]
The scalar parameter $R_0$ is positive. This inequality is used only
for analysis; the sampling algorithm below computes no matrix inverse.

\begin{lemma}[Sparse cumulative partition graphs]\label{lem:prefix-core}
Let $X$ have $N$ elements of positive integer masses $\mu_x$, with
$\sum_x\mu_x\le n$. Let $q\ge1$ be polynomially bounded. Given
partitions $Q_1,\ldots,Q_q$ and nonnegative coefficients $w_1,\ldots,w_q$
with the encodings specified above, define a graph at each index
$\tau\in\{1,\ldots,q+1\}$ by
\[
 k^\tau_{xy}=\mu_x\mu_y\sum_{i<\tau}w_i\join_{Q_i}(x,y).
\]
There is a sequential procedure that, while reading one partition at a time,
keeps a graph with capacities $\bar k^\tau$ and $\wt O(N)$ edges
satisfying, with high probability simultaneously for every such $\tau$
and every $C\subseteq X$,
\[
 k^\tau(\cut C)\le\bar k^\tau(\cut C)\le2k^\tau(\cut C).
\]
Its total work is $\wt O(qN)$ and its workspace is $\wt O(N)$.
It is combinatorial in the sense defined in Section~\ref{sec:definitions}.
\end{lemma}
\begin{proof}
Let $0<\eta<1/2$ and $0<\zeta<1/2$ be per-call error and failure
parameters; their values are specified below. Omit blocks with $w_i=0$
and blocks of size one. A remaining block $B$ contributes the clique
with capacity $w_i\mu_x\mu_y$ on $xy$. Put $R=\mu(B)$. The effective-resistance leverage
of edge $xy$ in this clique is $(\mu_x+\mu_y)/R$, and the sum of these
values is $|B|-1$. To verify the first formula, its Laplacian is
$w_i(R\diag(\mu)-\mu\mu^{\mathsf T})$, where $\mu$ is the column
vector $(\mu_x)_{x\in B}$ and $\diag(\mu)$ is the diagonal matrix of
those masses. The value $(e_x-e_y)^{\mathsf T}L^+(e_x-e_y)$ is
$(\mu_x^{-1}+\mu_y^{-1})/(w_iR)$. Here $e_x$ is the coordinate unit
vector in $\mathbb R^B$, the Laplacian $L$ of a graph has diagonal entries $\sum_yc_{xy}$
and off-diagonal entries $-c_{xy}$, and $L^+$ is its inverse on the
orthogonal complement of its kernel and zero on its kernel.
This quadratic form defines effective resistance. Leverage is capacity
times effective resistance.

Sample
\[
 s=3(|B|-1)\lceil\eta^{-2}\rceil
                 \left\lceil\log_2\frac{2(|B|-1)}{\zeta}\right\rceil
   =O(|B|\eta^{-2}\log(n/\zeta))
\]
independent edges with probability $p_{xy}=(\mu_x+\mu_y)/(R(|B|-1))$, assigning each sampled
copy capacity $w_i\mu_x\mu_y/(sp_{xy})$. This is the leverage sampler
of Spielman--Srivastava~\cite{SS11}. To obtain the stated failure
probability, put $\ell=|B|-1$, write $p_e=p_{xy}$ for $e=xy$, and put
\[
 v_e=\sqrt{w_i\mu_x\mu_y}\,L^{+1/2}(e_x-e_y).
\]
Here $L^{+1/2}$ is the positive semidefinite square root of $L^+$.
These vectors lie in the space orthogonal to the all-one vector,
which has dimension $\ell$. In any orthonormal coordinates on that
space,
\[
 \sum_e v_ev_e^{\mathsf T}=I_\ell,\qquad
 v_e^{\mathsf T}v_e=\ell p_e.
\]
The first identity follows by substituting
$L=\sum_{e=xy}w_i\mu_x\mu_y(e_x-e_y)(e_x-e_y)^{\mathsf T}$;
the second is the leverage formula above. If $e_j$ is sample $j$,
set $Y_j=v_{e_j}v_{e_j}^{\mathsf T}/(sp_{e_j})$. Then
$\sum_j\E Y_j=I_\ell$ and $0\preceq Y_j\preceq(\ell/s)I_\ell$,
the latter by Cauchy--Schwarz. The displayed matrix concentration
inequality bounds failure by $2\ell\exp(-\eta^2s/(3\ell))$.
The chosen sample size makes this at most $\zeta$.
Undoing multiplication by $L^{+1/2}$ gives a $1\pm\eta$
approximation to every Laplacian quadratic form. For a cut, use the
vector with entry one on its selected side and zero elsewhere;
its quadratic form is exactly cut capacity. Thus all cuts are
approximated simultaneously.

The implementation needs no linear-system computation: choose the first
endpoint proportionally to $\mu_x$, and the second uniformly among
the other vertices. The resulting unordered probability is exactly
$p_{xy}$. A block of size one contributes no edges. The vectors $v_e$
and the matrices $Y_j$ are used only in the proof.

Maintain a binary counter of consecutive intervals of processed layers.
Each counter entry is a sparsifier of its interval sum. Insert the current
layer sparsifier and, when two equally long intervals meet, add their
graphs and apply part~\ref{ext:bk}. At most
$d_0=1+\lceil\log_2 q\rceil$ errors accumulate in any entry, including
its initial clique sampling. Take $\eta=1/(10d_0)$ and scale the union
of current entries by $(1-\eta)^{-d_0}$. Each original cut capacity is
then preserved between factors one and
$((1+\eta)/(1-\eta))^{d_0}<2$.
There are fewer than $q$ merges, each on $\wt O(N)$ edges, and at most
$qN$ initial block samplings. For any desired failure exponent $a>0$,
use $\zeta=n^{-a}/(4nq)$ in each sampling and compression call.
A union bound proves the simultaneous assertion.
The probabilities in the explicit clique sampler are rational numbers
with polylogarithmic encodings. The stated compression theorem handles
capacitated graphs; rational encodings and inverse-polylogarithmic errors
contribute only polylogarithmic factors here.
\end{proof}

\section{A single graph for all distance thresholds}
\label{sec:sketch}

This section concerns a graph $G$ on $V$ given by a finite stream of distinct edge records with unit
capacities and nonnegative keys $\kappa_e$. Parallel records receive distinct
identifiers. Set $m_0$ to an upper bound on its number of records, and order
edges lexicographically by $(\kappa_e,\text{identifier})$. For a graph $R$ on $V$ and a real threshold $t$,
$R_{\le t}$ retains its edges with key at most $t$. A spanning forest contains
a spanning tree of each connected component; a minimum spanning forest
minimizes the sum of edge keys, with the stated tie-breaking order.
Define $C_k(R)$ by taking $k$ successive minimum spanning forests,
each in the graph remaining after the preceding forests are removed.

\begin{lemma}[Prefix certificates]\label{lem:msf}
For every $t,S$, if $q=|\cut_{R_{\le t}}S|$, then
\[
 |\cut_{C_k(R)_{\le t}}S|\ge\min\{k,q\}.
\]
\end{lemma}
\begin{proof}
Kruskal's algorithm processes edges in increasing order and accepts an edge
when it joins distinct forest components. Thus its restriction to a key
prefix is a spanning forest of that prefix. The same statement applies in
each residual graph. Whenever a residual prefix has an edge crossing $S$,
the spanning tree of its component has a crossing edge. Each successive
forest therefore removes at least one crossing edge until none remains.
\end{proof}

A \emph{$k$-cut certificate} is an edge subset whose intersection with
every cut has at least the smaller of $k$ and the original cut's edge
count. Thus Lemma~\ref{lem:msf} provides a $k$-cut certificate for
each prefix simultaneously.

An edge is \emph{$r$-weak} if its strength $k_e$, defined in
Theorem~\ref{thm:external}\ref{ext:bk}, is less than $r$.
We give the certificate fact used below; it is also Theorem 8.3 and
footnote 2 of~\cite{BK15}.

\begin{lemma}[Weak edges occur in the certificate]\label{lem:weak}
For an unweighted graph on $n$ vertices and $r\ge1$, the union of
$\lceil4r\rceil\lceil\log_2(n+1)\rceil$ successive spanning forests
contains every $r$-weak edge. The forests need not be minimum spanning
forests.
\end{lemma}
\begin{proof}
Call a vertex set with at least two vertices $r$-strong if its induced
subgraph has minimum cut at least $r$. A singleton is also called
$r$-strong. The union of two intersecting $r$-strong sets is
$r$-strong: any nontrivial cut of the union splits at least one of
the two sets. Thus maximal $r$-strong sets partition the vertices.
Contract these sets. Every $r$-weak edge remains uncontracted, and
the contracted graph has no $r$-strong set with at least two vertices.
Indeed, the union of the contracted sets in such a set would be
$r$-strong. A cut splitting a contracted set has capacity at least
$r$ inside that set; a cut splitting none has the corresponding
contracted-graph capacity at least $r$. This contradicts maximality.

Any unweighted graph on $s\ge2$ vertices with no nontrivial $r$-strong
induced subgraph has fewer than $r(s-1)$ edges. To prove this, repeatedly
split a nonsingleton vertex set along a cut of capacity less than $r$.
There are $s-1$ splits, each edge is counted at its first separating
split, and all edges are eventually counted.
Apply this observation to the vertices incident to an edge of the
contracted graph; if there are none, the proof is complete.
Their average degree is less than $2r$, so more than half have degree
less than $4r$. In one block of $\lceil4r\rceil$ forest removals,
every edge incident to any such vertex is removed: the forests form
a cut certificate by the proof of Lemma~\ref{lem:msf}, and the
corresponding singleton cut in the contracted graph has fewer than $4r$ edges.
Pulling that cut back to the uncontracted graph proves the assertion
even when edges inside a contracted set are present. The remaining
contracted graph still has no nontrivial $r$-strong induced subgraph. Hence
each block halves its number of nonisolated vertices. After
$\lceil\log_2(n+1)\rceil$ blocks no edge of the contracted graph remains.
In particular every $r$-weak edge has been removed.
\end{proof}

By Lemma~\ref{lem:msf}, restriction of successive minimum spanning
forests to a key prefix gives successive spanning forests of that
prefix. Lemma~\ref{lem:weak} therefore applies to every prefix.

\begin{lemma}[One halving step]\label{lem:halving}
Let $0<\eta<1/2$ and let $0<\zeta<1/2$. For
\[
 r=C\eta^{-2}\log(nm_0/\zeta),\qquad
 k=\lceil8r\rceil\lceil\log_2(n+1)\rceil,
\]
keep $C_k(R)$, and retain each other edge independently with probability
one half and capacity two. The resulting graph approximates every cut of
every key prefix within $1\pm\eta$, except with probability $\zeta$.
\end{lemma}
\begin{proof}
Fix a prefix. Its certificate contains every edge of strength below $2r$,
by Lemma~\ref{lem:weak}. Every unprotected edge has strength at
least $2r$, so its sampling probability one half is at least $r/k_e$.
Protected edges have sampling probability one. Apply
Theorem~\ref{thm:external}\ref{ext:bk}, with failure probability
$\zeta/m_0$ for this prefix. There are at most $m_0$ nonempty distinct key
prefixes. A union bound proves the result.
\end{proof}

\begin{theorem}[One-pass prefix cut sparsifier]\label{thm:prefix}
For $0<\varepsilon\le1/4$, $0<\delta<1/2$, and $m_0=n^{O(1)}$, there is
a combinatorial one-pass algorithm retaining a keyed capacitated graph $H$
with
\[
 O\!\left(n\varepsilon^{-2}\log^5(nm_0/\delta)\right)
\]
edges, such that, with probability at least $1-\delta$,
\begin{equation}\label{eq:prefix}
 (1-\varepsilon)|\cut_{G_{\le t}}S|
 \le a_{H_{\le t}}(\cut S)
 \le(1+\varepsilon)|\cut_{G_{\le t}}S|
 \quad\text{for every }t,S.
\end{equation}
Each retained edge has its original key and capacity $2^j$ for a
nonnegative integer $j=O(\log(m_0/\delta))$. For fixed accuracy and
inverse-polynomial failure, updates take $\wt O(1)$ time per record
and space is $\wt O(n)$.
\end{theorem}
\begin{proof}
Put $L=\lceil\log_2(3m_0/\delta)\rceil$ and
$\eta=\varepsilon/(4L)$.
Starting with $R_0=G$, at level $j$ protect $C_k(R_j)$ and put its
edges in $H$ with capacity $2^j$. Independently retain each unprotected
edge with probability one half to form $R_{j+1}$. Use
Lemma~\ref{lem:halving} at every level with failure $\delta/(3L)$.
Each substitution $R_j\mapsto C_k(R_j)+2R_{j+1}$ has relative error at
most $\eta$. Substitution from the last level backwards bounds the
composed relative factors by $(1\pm\eta)^L$, which lie in
$[1-\varepsilon,1+\varepsilon]$. An edge in $R_L$ has survived $L$
independent halving decisions; thus
$\Prb[R_L\ne\varnothing]\le m_02^{-L}\le\delta/3$.
If $R_L$ is nonempty, report failure. There are at most $k(n-1)L$
protected edges, giving the size bound.

Here is the one-pass implementation. Maintain all $k$ forests at every
level. Given one insertion $e$ into a minimum spanning forest $F$ of a graph
$A$, add $e$ if it joins two components. Otherwise let $f$ be the
largest-key edge on its forest path. Replace $f$ by $e$ if $e<f$,
and emit $f$; if $f<e$, emit $e$ without changing the forest.
If no edge is emitted, stop this insertion. The cycle rule proves that
the new forest is $F'=\MSF(A+e)$ and that
\[
 (A+e)\setminus F'=(A\setminus F)+o,
\]
where $o$ is the emitted edge, or the empty set if no edge is emitted.
Therefore the residual only receives insertions. Feed the emitted edge
into the next forest; after the $k$th forest, toss its halving coin and,
if retained, feed it into the first forest at the next level.
Induction on insertions and forest indices proves that the final forests
are exactly the certificates of the conceptual construction.
Coins at one level do not affect its own certificate. Conditional on the
graph entering that level, its residual coins are independent, as required
by Lemma~\ref{lem:halving}. An edge cannot leave and subsequently reenter
the same forest cascade.

Set a failure flag when an edge reaches $R_L$ and discard that edge.
This enforces the size bound on all executions. Link, cut, and
maximum-on-path operations are supported in $O(\log n)$ amortized time
by dynamic trees~\cite{ST83}. One original record traverses at most $kL$
forests. This proves the resource bounds.
\end{proof}

All subsequent uses of a failed construction return the height-$M$
ultrametric and count the failure against their stated probability. Every intermediate sparsifier is also given its asserted $C N\log^c(n+1)$ edge cap, for sufficiently large fixed constants $C,c$; exceeding a cap marks that trial failed. Thus the working-space bounds hold on every execution, including failed executions.

\begin{lemma}[A stored graph evaluates $L_1$]\label{lem:cost-sketch}
For the complete input, maintain also $D_0=\sum_{xy}d_{xy}$.
For any ultrametric $u$ with entries in $[0,M]$, define
\[
 \widehat L_H(u)=D_0-\sum_{xy\in\binom V2}u_{xy}
       +2\sum_{xy\in E(H)}a_{xy}(u_{xy}-d_{xy})_+.
\]
On event~\eqref{eq:prefix}, simultaneously for every such $u$,
\[
 |\widehat L_H(u)-L_1(d,u)|\le2\varepsilon L_1(d,u).
\]
Thus selecting a candidate of least $\widehat L_H$ costs at most
$(1+2\varepsilon)/(1-2\varepsilon)$ times the least true candidate cost.
Evaluation uses $\wt O(n)$ time per hierarchy.
\end{lemma}
\begin{proof}
For a partition $P$, let $p_t(P)$ be the number of edges of $G_{\le t}$
between its blocks. Summing the cut inequalities over its blocks and
dividing by two preserves $p_t(P)$ within $1\pm\varepsilon$.
For $P=P_u(t)$, integration gives
\[
 \int_0^M p_t(P_u(t))\,dt=\sum_{xy}(u_{xy}-d_{xy})_+.
\]
Its stored-graph version is the last sum in $\widehat L_H$, before
multiplication by two. The identity
$|d-u|=d-u+2(u-d)_+$ proves the bound and candidate-selection assertion.
If $h_v$ is an internal tree height and its children have leaf counts
$n_1,\ldots,n_j$, the pairs with lowest common ancestor $v$ number
$\sum_{i<i'}n_in_{i'}$. Summing $h_v$ times this number computes
$\sum_{xy}u_{xy}$ in linear tree time. Queries on the $\wt O(n)$ stored
edges can use ancestor searches in $O(\log n)$ time each, with linear-space
tree preprocessing, or an Euler-tour lowest-common-ancestor data structure.
\end{proof}

\section{A common partition sequence}
\label{sec:guide}

For $t\in[0,M)$ let $G_t$ have positive pairs $xy$ with $d_{xy}\le t$;
all other pairs are negative. For a partition $P$, its
\emph{disagreement cost} at $t$ is
\[
 \operatorname{cc}_t(P)=
 \sum_{xy\in\binom V2}
 |\ind[d_{xy}\le t]-\join_P(x,y)|.
\]
For partition-valued functions $P,Q$ on $[0,M)$, put
\[
 D(P,Q)=\int_0^M\sum_{xy}|\join_{P(t)}(x,y)-\join_{Q(t)}(x,y)|\,dt,
 \qquad
 D(G,P)=\int_0^M\operatorname{cc}_t(P(t))\,dt.
\]
These functions will always be piecewise constant with finitely many
intervals, so the integrals are finite sums. For every ultrametric $u$,
\begin{equation}\label{eq:cake}
 L_1(d,u)=D(G,P_u),
\end{equation}
because the disagreement interval for one pair has length $|d_{xy}-u_{xy}|$.

Choose a uniform random permutation $\pi$ of $V$; $\pi(x)$ is its rank.
At each threshold define
\begin{equation}\label{eq:guide}
 f_t(v)=\arg\min_{x\in\{v\}\cup\{y:d_{vy}\le t\}}\pi(x),\qquad
 g_t(v)=
 \begin{cases}f_t(v),&f_t(f_t(v))=f_t(v),\\
 v,&f_t(f_t(v))\ne f_t(v).
 \end{cases}
\end{equation}
The \emph{guide partition} $Q(t)$ groups vertices with equal labels $g_t$.
A vertex $r$ with $f_t(r)=r$ is a \emph{center}. A non-singleton block has
a center as its label and consists of positive neighbors of that center.

\begin{lemma}[Guide cost]\label{lem:guide}
For every fixed partition $P$ at one threshold,
$\E_\pi\operatorname{cc}_t(Q(t))\le7\operatorname{cc}_t(P)$.
Consequently, for every fixed ultrametric $u$,
\[
 \E D(G,Q)\le7L_1(d,u),\qquad
 \E D(Q,P_u)\le8L_1(d,u).
\]
\end{lemma}
\begin{proof}
Fix a block $C$ of $P$, let $k=|C|$, and let $r_C$ be its minimum-ranked
vertex. Let $z_C$ be its number of internal negative pairs. For $v\in C$
let $A_v$ be its number of positive neighbors outside $C$, and let
$A_C=\sum_{v\in C}A_v$.
If $r_C$ is a center, let $B_C=\{v\in C:g_t(v)\ne r_C\}$; otherwise
put $B_C=C$. The number of positive pairs of $C$ separated by the guide is
at most $k|B_C|$.

When $r_C$ is a center, a vertex in $B_C$ either is negatively adjacent
to $r_C$, or has an outside positive neighbor preceding all of $C$.
Their expected numbers are at most $2z_C/k$ and $A_C/(k+1)$.
The center test fails only if an outside positive neighbor of $r_C$
precedes $r_C$. For a fixed positive pair $vw$, with $v\in C,w\notin C$,
the event that $w$ precedes all of $C$ and $v=r_C$ has probability
$1/(k(k+1))$. Thus
\[
 \E[k|B_C|]\le2z_C+2A_C.
\]
If the comparison partition has $p$ positive crossing pairs and $z$ negative
internal pairs, summing gives at most $2z+4p$ new separated positive pairs.

For a center $r$, let $S_r$ be its guide block, $C(r)$ its comparison
block, and $L_r=|S_r\setminus C(r)|$. New negative pairs joined across
comparison blocks are at most $|S_r|L_r$. Deterministically
$|S_r|\le|N^+[r]|$ and $L_r\le A_r$, where
$N^+[r]=\{r\}\cup\{v:rv\text{ positive}\}$. The center probability is
$1/|N^+[r]|$. Hence the expected sum of these products is at most
$\sum_rA_r=2p$. Previously disagreeing pairs contribute at most $p+z$.
The total is at most $7p+3z\le7(p+z)$.
Apply this at each fixed $t$ with $P=P_u(t)$, integrate, and use the
triangle inequality for $D$.
\end{proof}

\begin{lemma}[Exact small representation]\label{lem:records}
The functions $f_t,g_t$ can be stored exactly, with high probability, in
$\wt O(n)$ words, maintained in one pass with $\wt O(n^2)$ total update
time. They have a common grid of $\wt O(n)$ interval endpoints.
Each label query takes $\wt O(1)$ time.
\end{lemma}
\begin{proof}
For each vertex store an initial record $(-\infty,v)$ and, in increasing
distance order, one record whenever its minimum neighbor rank strictly
decreases. Batch equal keys. The last record with key at most $t$ gives
$f_t(v)$. Two such queries give $g_t(v)$.

At a fixed input prefix, order the available neighbors by distance and
fixed identifier, placing $v$ first. Their relative ranks are uniformly
random. The number $R_m$ of left-to-right minimum ranks among $m\le n$
vertices has the same distribution as $\sum_{j=1}^m B_j$, where the
$B_j$ are independent, take values in $\{0,1\}$, and satisfy
$\Prb[B_j=1]=1/j$. To see this, the rank of the $j$th vertex among the
first $j$ vertices is uniform on $\{1,\ldots,j\}$. These insertion
ranks are independent, since their possible sequences are in bijection
with permutations. A new minimum occurs exactly when that rank is one. Therefore
\[
 \E e^{R_m}=\prod_{j=1}^m(1+(e-1)/j)
 \le \exp((e-1)(1+\log n)).
\]
Exponential Markov inequality gives $R_m\le C\log(n+1)$ with any required
inverse-polynomial failure. Union-bound over vertices, input prefixes,
and repetitions. Equal-key batching only reduces the count.
Enforce this bound as a deterministic cap; an overflow marks the trial
failed.

For an incoming pair $vx$ of distance $z$, query the current minimum at
$z$. Discard $x$ if its rank is no smaller. Otherwise insert $(z,x)$,
replace an equal-key record if present, and delete later records whose
ranks are no smaller than $\pi(x)$, stopping at the first smaller rank.
This is exactly the minimum function of the candidates received so far.
A discarded candidate cannot become a minimum after further insertions.
Every inserted record is deleted at most once. Store the capped lists as
sorted arrays. A query, insertion, or deletion scan costs $O(\log n)$
word operations because a list has $O(\log n)$ entries. This gives the
update bound. Both endpoints are updated.

The union of final record keys, together with $0,M$, is a common grid:
$f_t(v)$ changes only at a key of its own list, and $f_t(f_t(v))$ changes
only at a key of one of the lists. Thus $g_t(v)$ also changes only on this
grid. The total number of keys is $\wt O(n)$. Alternatively, sorting all
input pairs and maintaining the same minima constructs these lists
offline in $\wt O(n^2)$ time.
\end{proof}

Fix any grid $0=\theta_0<\theta_1<\cdots<\theta_q=M$ on whose open
intervals $Q$ is constant, and put $w_i=\theta_i-\theta_{i-1}$ and
$Q_i=Q(\theta_{i-1})$. Threshold endpoints have measure zero.
We will use either the grid of all input distances, of size $O(n^2)$,
or the record grid of Lemma~\ref{lem:records}.
Adjoin a discrete partition at index zero and an indiscrete partition
at index $q+1$, both with weight zero. Here \emph{discrete} means all
singleton blocks and \emph{indiscrete} means the single block $V$.
Thus $w_0=w_{q+1}=0$. Every subsequent sum over an unspecified layer
index $i$ ranges over $i=0,\ldots,q+1$.

\begin{lemma}[Grid restriction]\label{lem:grid}
For every nested partition function $P$, there is a nested function
$\widehat P$ changing only at grid endpoints such that
$D(Q,\widehat P)\le D(Q,P)$. A nested sequence on this grid specifies
an ultrametric with first-join height $\theta_{i-1}$ when a pair first
joins at index $i$; a first join at $q+1$ has height $M$.
\end{lemma}
\begin{proof}
Within one interval, $Q$ is fixed. The successive partitions of $P$ form
a chain. Its disagreement integral is a linear function of the nonnegative
lengths spent in its states, with their sum fixed. Assign the whole interval
to a minimum-cost state, pushing preceding transitions to the left endpoint
and succeeding transitions to the right endpoint. This preserves the order
of transitions and does not increase cost. Repeat on all intervals.
First-join heights satisfy the ultrametric inequality because every
threshold relation is an equivalence relation.
\end{proof}

\section{Weighted partition states and projection}
\label{sec:projection}

A \emph{partition state} consists of a set $X$ of $N$ elements called
\emph{atoms}, integer masses $\mu_x\ge1$, release indices
$\lambda_x\in\{1,\ldots,q+1\}$, an upper index $H\in\{1,\ldots,q+1\}$,
and partitions $Q_i$ of $X$. Write $\mu(C)=\sum_{x\in C}\mu_x$ and
assume $\mu(X)\le n$ and $\lambda_x\le H$.
The partitions have $x$ singleton for $i<\lambda_x$ and are indiscrete
for $i\ge H$; the guide need not be nested.
A \emph{feasible comparator} is a nested sequence with those same
boundary properties. For any two partition sequences set
\[
 D_\mu(Q,P)=\sum_iw_i\sum_{xy\in\binom X2}
       \mu_x\mu_y|\join_{Q_i}(x,y)-\join_{P_i}(x,y)|.
\]
At the root, $X=V$, all masses are one, all releases are one, and $H=q+1$.
In the algorithms, each atom represents a nonempty subset of $V$ and
its mass is the cardinality of that subset. Within one state these
subsets are disjoint. Contracting a set of atoms replaces their subsets
by their union. Thus an atom pair $xy$ represents exactly $\mu_x\mu_y$
original pairs. Pairs within a contracted atom are handled in the
inside child and are absent from the outside child's objective.

For $\varnothing\ne C\subsetneq X$, $a\in C$, and
$\max_{x\in C}\lambda_x\le\tau\le H$, define on a partition $\Pi$
\[
 S_C\Pi=\{B\cap C,B\setminus C:B\in\Pi\}\setminus\{\varnothing\},
\]
\[
 M_{C,a}\Pi=\{C\cup(\Pi(a)\setminus C)\}
 \cup\{B\setminus C:B\in\Pi,\ B\ne\Pi(a)\}\setminus\{\varnothing\}.
\]
The \emph{projection} $T_{C,\tau,a}$, abbreviated $T_a$ when $C,\tau$
are fixed, applies $S_C$ at $i<\tau$ and $M_{C,a}$ at $i\ge\tau$.
We call $\tau$ the \emph{birth index} of $C$. A \emph{separator} is a
pair $(C,\tau)$ with $2\le |C|\le N-1$ and
$\max_{x\in C}\lambda_x\le\tau\le H$. It is \emph{feasible} when
these inequalities hold. Separators are used only when $N\ge3$.
A set $C$ is \emph{compatible} with a nested sequence $P$ at $\tau$
if no block crosses $C$ below $\tau$ and all of $C$ lies in one block
at and above $\tau$.

\begin{lemma}[Projection identities]\label{lem:projection}
Projection maps a nested sequence to a nested sequence and preserves
the boundary properties. A compatible
comparator is fixed. If $A\in C$ is sampled with probability
$\Prb[A=x]=\mu_x/\mu(C)$, then, for arbitrary partition sequences $Q,P$,
\[
 \E_A D_\mu(T_AQ,T_AP)\le D_\mu(Q,P).
\]
\end{lemma}
\begin{proof}
Splitting along a fixed set preserves refinement. In a nested sequence
the anchor blocks grow, so the coarse merging also preserves refinement.
Across $\tau$, each fine block lies in the corresponding coarse block.
Splitting cannot merge unreleased atoms. A new coarse merge involves
released atoms of $C$; an unreleased outside atom was a singleton and
cannot lie in the old anchor block. The upper boundary is unchanged.
Compatibility makes each relation unchanged.

At a fine layer, crossing disagreements are deleted and other
disagreements are unchanged. At a coarse layer, internal disagreements
are deleted and outside disagreements are unchanged. For $y\notin C$,
the total new crossing disagreement is
$\mu(C)\mu_y|\join_{Q_i}(A,y)-\join_{P_i}(A,y)|$. Its expectation is exactly
$\sum_{x\in C}\mu_x\mu_y|\join_{Q_i}(x,y)-\join_{P_i}(x,y)|$.
Sum over $y$ and indices.
\end{proof}

For one layer $\Pi$, put $c=\mu(C)$,
$p_B=\mu(B\cap C)$, $q_B=\mu(B\setminus C)$, $r_B=\mu(B)$, and
$B_a=\Pi(a)$. Define
\[
 J_\Pi(C)=\sum_{\{B,B'\}\in\binom\Pi2}p_Bp_{B'},\qquad
 s_\Pi(C)=\sum_Bp_Bq_B,
\]
\[
 h_{\Pi,a}(C)=J_\Pi(C)+\sum_{B\ne B_a}p_Bq_B+q_{B_a}(c-p_{B_a}),
\]
\[
 \widetilde h_{\Pi,a}(C)=J_\Pi(C)+\sum_{B\ne B_a}p_Bq_B+r_{B_a}(c-p_{B_a}),
\]
and
\[
 \Gamma_{\tau,a}(C)=\sum_{i<\tau}w_is_{Q_i}(C)
             +\sum_{i\ge\tau}w_i\widetilde h_{Q_i,a}(C).
\]
The symbol $\Gamma$ always denotes this certificate.

\begin{lemma}[Certificate bounds]\label{lem:certificate}
Fine projection edit is $s_\Pi(C)$ and coarse projection edit is
$h_{\Pi,a}(C)$. For every fixed $a\in C$ and an independent
mass-proportional $A\in C$,
\[
 \E_A D_\mu(Q,T_AQ)\le2\Gamma_{\tau,a}(C).
\]
For a compatible feasible comparator $P$,
\[
 \E_{a\sim\mu|_C}\Gamma_{\tau,a}(C)\le4D_\mu(Q,P).
\]
The notation $a\sim\mu|_C$ means $\Prb[a=x]=\mu_x/c$ for $x\in C$.
\end{lemma}
\begin{proof}
A split separates precisely the old joined crossing pairs.
A merge joins internal pairs from distinct old blocks, separates old
crossing pairs in nonanchor blocks, and joins the outside anchor-block
points to the included points in other blocks. These give the formulas.
Averaging coarse edit gives
\[
 J_\Pi(C)+2\sum_Bp_Bq_B(1-p_B/c).
\]
For $B\ne B_a$, the summand before multiplication by two is at most
$p_Bq_B$; for $B=B_a$ it is at most $q_{B_a}(c-p_{B_a})$.
This proves the first inequality.
If $P$ is compatible, $T_AP=P$, so the triangle inequality and
Lemma~\ref{lem:projection} bound expected exact edit by $2D_\mu(Q,P)$.
The excess of $\widetilde h$ over $h$, averaged over $a$, is
\[
 \frac1c\sum_Bp_B^2(c-p_B)
 \le\sum_Bp_B(c-p_B)=2J_\Pi(C).
\]
At coarse layers all pairs counted by $J_\Pi(C)$ are comparator
disagreements. Their weighted sum is at most $D_\mu(Q,P)$.
\end{proof}

\begin{lemma}[Compatible balanced sets]\label{lem:balanced}
Every feasible comparator with $N>6$ has a compatible set $C$ at a feasible
$\tau$ with $N/3\le|C|\le2N/3$. For $3\le N\le6$, a proper compatible set
with $2\le|C|\le N-1$ exists.
\end{lemma}
\begin{proof}
Use its hierarchy tree with atoms as leaves. Descend to a vertex whose
children each have at most $2N/3$ leaves while the current subtree has
more than $2N/3$. If a child has at least $N/3$ leaves, take it.
Otherwise accumulate children until their union first reaches $N/3$;
the union has fewer than $2N/3$ leaves. Insert their union at their
parent's height. This is an equal-height refinement and changes no pair
distance. Releases are no later than that height. For the small cases,
take a proper internal subtree with at least two leaves if one exists.
If none exists, every root child is a leaf after one-child vertices are
suppressed, and a union of two root children has between two and $N-1$
leaves. Insert that union at the root height.
\end{proof}

\begin{lemma}[Sizes in a balanced recursion]\label{lem:recursion-size}
Consider a recursion starting at count $n$. Each state of count $N$
where recursion continues has child counts $k,N-k+1$, with
$2\le k\le N-1$.
Suppose there are constants $\rho<1$ and $N_0$ such that both child
counts are at most $\rho N$ whenever $N>N_0$, and no state of count
at most $N_0$ has more than $N_0$ further levels. Then the number of
states is $O(n)$, the depth is $O(\log n)$, and
\[
 \sum_v N_v=O(n\log n),\qquad
 \sum_v N_v^s=O(n^s)\quad\text{for every fixed }s>1.
\]
Here $v$ ranges over states and $N_v$ is a state's count.
\end{lemma}
\begin{proof}
The quantity $N-1$ is positive for every state of count at least two,
and its child values sum to $N-1$. Thus there are at most $n-1$
states where recursion stops, and at most $2n-3$ states in
their binary recursion tree. Every child count is at least two.
The depth assertion follows
from the factor $\rho$ above count $N_0$.
At any fixed depth, the sum of $N_v-1$ is at most $n-1$, so the sum
of $N_v$ is at most $2(n-1)$. This proves the first sum bound.
For states with $N_v>N_0$ at depth $j$, every ancestor is also above
$N_0$, whence $N_v\le\rho^j n$. Their contribution to the second sum
at that depth is at most $2(n-1)(\rho^j n)^{s-1}$.
The sum over $j$ is a convergent geometric series times $O(n^s)$.
There are $O(n)$ states of count at most $N_0$, each contributing a
constant depending only on $s,N_0$. This proves the second bound.
\end{proof}

\begin{lemma}[Recursion principle]\label{lem:l1-recursion}
Suppose a procedure chooses a feasible $C,\tau$ with
$2\le|C|\le N-1$, constant-factor count balance when $N>16$, and
\[
 \E D_\mu(Q,T_AQ)\le K D_\mu(Q,P)
\]
for every feasible comparator $P$, using a fresh mass-proportional
execution anchor $A$. Project, recurse inside $C$ with upper index $\tau$,
and recurse outside with $C$ contracted to one atom of mass $\mu(C)$
and release $\tau$. Substitute the inside hierarchy for that atom.
Then its output $R$ satisfies
\[
 \E D_\mu(Q,R)\le(1+Kh(N))D_\mu(Q,P),
\]
where $h(N)=O(\log n)$ is the maximum nontrivial recursion depth.
\end{lemma}
\begin{proof}
Project $P$ using the same anchor. Its inside and outside restrictions
are feasible. The two comparator distances sum exactly to
$D_\mu(T_AQ,T_AP)$: internal pairs occur inside, outside pairs outside,
and crossing pairs have one common relation and aggregate weight
$\mu(C)\mu_y$. The same identity holds for substituted outputs.
Substitution preserves nesting because inside merges occur no later
than $\tau$ and outside merges involving the contracted atom occur no
earlier than $\tau$.
Explicitly, below $\tau$ the inside guide is $Q_i$ restricted to $C$,
and at and above $\tau$ it is the single block $C$. Below $\tau$ the
outside guide is $Q_i$ restricted to $X\setminus C$, together with
the singleton contracted atom. At and above $\tau$, start with the
same restriction and insert that atom into $Q_i(A)\setminus C$,
creating a singleton if this set is empty.
The same rules apply to the projected comparator. These definitions
specify both recursive inputs without retaining the projected sequence.
For the induction, put $B=D_\mu(Q,P)$ and let $h$ be the remaining
depth bound. Conditional on the separator and its execution anchor,
the two inductive error bounds sum to at most
$(1+K(h-1))D_\mu(T_AQ,T_AP)$. The triangle inequality and
Lemma~\ref{lem:projection} therefore give
\[
 \E D_\mu(Q,R)
 \le KB+(1+K(h-1))B=(1+Kh)B.
\]
The two child counts are $|C|$ and $N-|C|+1$. Constant balance gives
geometric decrease above the fixed small-state size; below it both
counts decrease. Lemma~\ref{lem:recursion-size} gives the depth bound,
completing the induction.
At $N=2$ choose the best feasible merge index
by a scan; at $N=1$ return the single leaf.
\end{proof}



\section{The combinatorial offline \texorpdfstring{$L_1$}{L1} algorithm}
\label{sec:slow}

For a partition state define, for $x\ne y$,
\[
 b^\tau_{xy}=\sum_{i\ge\tau}w_i\ind[Q_i(x)\ne Q_i(y)],
 \qquad B_\tau(C)=\sum_{xy\in\binom C2}\mu_x\mu_yb^\tau_{xy}.
\]
Set $b^\tau_{xx}=0$. Each $b^\tau$ is a pseudometric, meaning a symmetric
nonnegative function with zero diagonal satisfying the triangle inequality:
each summand is the distance zero within a partition block and one between
blocks. For a certificate anchor $a$, put
\[
 \kappa^{\tau,a}_{xy}=\mu_x\mu_y\left(
  \sum_{i<\tau}w_i\ind[Q_i(x)=Q_i(y)]
 +\sum_{i\ge\tau}w_i\ind[Q_i(x)=Q_i(y)\ne Q_i(a)]\right),
\]
\[
 \ell^{\tau,a}_x=
 \mu_x\sum_{i\ge\tau}w_i\mu(Q_i(a))\ind[Q_i(x)\ne Q_i(a)].
\]
For any edge coefficients $v_{xy}$, write
$v(\cut C)=\sum_{x\in C,y\notin C}v_{xy}$; for unary coefficients
$v_x$, write $v(C)=\sum_{x\in C}v_x$.

\begin{lemma}[Cut representation and one extra anchor]\label{lem:medoid}
For $a\in C$,
\[
 \Gamma_{\tau,a}(C)=\kappa^{\tau,a}(\cut C)+\ell^{\tau,a}(C)+B_\tau(C).
\]
Writing $c=\mu(C)$, there is $r\in C$ with
\[
 B_\tau(C)\le c\sum_{x\in C}\mu_xb^\tau_{rx}
       \le2B_\tau(C).
\]
The first inequality holds for every $r\in C$.
\end{lemma}
\begin{proof}
At a fine layer the cut coefficient counts $s_\Pi(C)$. At a coarse
layer the three summands count, respectively,
$\sum_{B\ne B_a}p_Bq_B$, $r_{B_a}(c-p_{B_a})$, and $J_\Pi(C)$.
For the lower inequality sum
$b^\tau_{xy}\le b^\tau_{xr}+b^\tau_{ry}$ over internal unordered pairs.
For the upper inequality sample $r$ mass-proportionally; the expectation
of $c\sum_x\mu_xb^\tau_{rx}$ is exactly $2B_\tau(C)$.
\end{proof}

For $N>6$, enumerate every $(\tau,a,r,c)$ with $1\le\tau\le H$,
$a,r\in X$, and integer $1\le c\le n$. The graph on $X\cup\{t\}$
has pair capacities $\kappa^{\tau,a}_{xy}$ and terminal capacities
$\ell^{\tau,a}_x+c\mu_xb^\tau_{rx}$ on $xt$.
Require $a,r$ inside, $t$ and all $\lambda_x>\tau$ outside, exact selected
mass $c$, and counts between $\lceil N/3\rceil$ and $\lfloor2N/3\rfloor$.
Use a sampled tree and the two-index dynamic program of
Lemma~\ref{lem:tree-dp}. Skip infeasible choices. Among returned sets choose
the one of least exactly computed $\Gamma_{\tau,a}(C)$.
Sample an independent mass-proportional execution anchor and use
Lemma~\ref{lem:l1-recursion}.

\begin{lemma}[Offline separator]\label{lem:slow-separator}
For every feasible comparator $P$, the chosen separator has
\[
 \E D_\mu(Q,T_AQ)\le16\beta(n+1)D_\mu(Q,P).
\]
\end{lemma}
\begin{proof}
Let $C^*,\tau^*$ be the set of Lemma~\ref{lem:balanced} and put
$B=D_\mu(Q,P)$. Lemma~\ref{lem:certificate} gives an anchor $a^*$ with
$\Gamma_{\tau^*,a^*}(C^*)\le4B$. Choose $r^*$ as in
Lemma~\ref{lem:medoid} and $c^*=\mu(C^*)$.
At that enumerated tuple the graph objective dominates $\Gamma$ for
every feasible set and is at most $2\Gamma(C^*)\le8B$ on the comparator.
Tree domination and expected stretch bound the expected returned
graph cost by $8\beta(n+1)B$. Selection does not increase its certificate,
and the independent execution anchor costs at most twice that certificate.
\end{proof}

For $3\le N\le6$, enumerate proper non-singleton $C$, feasible $\tau$,
and $a\in C$, selecting least $\Gamma$. This costs $O(q)$ time for each
constant-size state: for each $C,a$, the birth scores are one prefix
sum of fine terms plus one suffix sum of coarse terms.
A compatible candidate and Lemma~\ref{lem:certificate} give expected edit
at most $8D_\mu(Q,P)$. For $N=2$ scan feasible merge indices and minimize
the single pair's disagreement; for $N=1$ return its leaf.

\begin{theorem}[Combinatorial offline $L_1$]\label{thm:slow}
There is a combinatorial $O(\log^2 n)$-approximation for $L_1$ in
$\wt O(n^{10})$ time.
\end{theorem}
\begin{proof}
Use the grid of all distinct input distances, so $q=O(n^2)$, and construct
the guide of Section~\ref{sec:guide}. In this implementation store its
partition labels explicitly. A direct scan of all input pairs at each
layer constructs them in $O(qn^2)$ time and $O(qn)$ words, without
record caps. By Lemmas~\ref{lem:slow-separator}
and~\ref{lem:l1-recursion}, expected distance to any fixed nested
comparator is multiplied by $O(\beta(n)\log n)=O(\log^2 n)$.
Apply Lemmas~\ref{lem:guide} and~\ref{lem:grid}, and then the triangle
inequality with the input relation in~\eqref{eq:cake}. This gives expected
input cost $O(\log^2 n)\OPT_1$. Markov's inequality gives constant success;
$O(\log n)$ independent trials and exact input-cost selection amplify it.
For this expectation argument, use an ideal execution with correct
external tree distributions. Each external failure is given an
inverse-polynomial probability budget. The implemented execution agrees
with the ideal execution until such a failure, so the sum of those
budgets bounds the additional failure probability.

For time, a state has $O(qnN^2)$ tuples. One graph has $O(N^2)$ edges.
Its tree construction costs $\wt O(N^5)$, its two-index dynamic program
costs $O(n^2N^3)$, and direct capacity computation costs $O(qN^2)$.
For every fixed $s>1$, Lemma~\ref{lem:recursion-size} gives
$\sum_vN_v^s=O(n^s)$, where $v$ ranges over states. Total time is therefore
\[
 \wt O\!\left(qn\sum_v(N_v^7+n^2N_v^5+qN_v^4)\right)
 =\wt O(qn^8+q^2n^5)=\wt O(n^{10}).
\]
The small-state scans cost $O(qn)$ and cannot change the exponent.
\end{proof}

\section{The fast \texorpdfstring{$L_1$}{L1} separator}
\label{sec:fast-l1}

Define a second certificate, denoted only by $\mathcal H$, by
\[
 H_{\Pi,a}(C)=J_\Pi(C)+\sum_{B\ne B_a}p_Br_B+r_{B_a}(c-p_{B_a}),
 \qquad
 \mathcal H_{\tau,a}(C)=\sum_{i<\tau}w_is_{Q_i}(C)
                   +\sum_{i\ge\tau}w_iH_{Q_i,a}(C).
\]
Let
\[
 k^\tau_{xy}=\mu_x\mu_y\sum_{i<\tau}w_i\join_{Q_i}(x,y),\quad
 u^{\tau,a}_x=\mu_x\sum_{i\ge\tau}w_i
  \bigl(\mu(Q_i(x))+\mu(Q_i(a))\bigr)\ind[Q_i(x)\ne Q_i(a)],
\]
\[
 Z^\tau_a(C)=\sum_{x\in C}\mu_xb^\tau_{ax},\qquad
 F_{\tau,a,\gamma}(C)=k^\tau(\cut C)+u^{\tau,a}(C)+\gamma Z^\tau_a(C).
\]
Here $\gamma>0$ is an algorithm parameter.

\begin{lemma}[The second certificate]\label{lem:unary}
For all $C\ni a$, $\Gamma_{\tau,a}(C)\le\mathcal H_{\tau,a}(C)$ and
\[
 \mathcal H_{\tau,a}(C)=k^\tau(\cut C)+u^{\tau,a}(C)+B_\tau(C).
\]
For a compatible comparator $P$,
$\E_{a\sim\mu|_C}\mathcal H_{\tau,a}(C)\le6D_\mu(Q,P)$.
Moreover $F_{\tau,a,\gamma}(C)\ge\mathcal H_{\tau,a}(C)$ if
$\gamma\ge\mu(C)$.
\end{lemma}
\begin{proof}
The first assertion replaces $q_B$ by $r_B\ge q_B$ in nonanchor blocks.
The displayed representation follows by counting each block contribution.
Averaging over $a$ gives the coarse value
$J_\Pi(C)+2\sum_Bp_Br_B(1-p_B/c)$.
Its excess over expected exact edit is
$2\sum_Bp_B^2(1-p_B/c)\le4J_\Pi(C)$.
Expected exact edit is at most $2D_\mu(Q,P)$ by
Lemma~\ref{lem:projection}; the weighted internal sum of $J$ is at most
$D_\mu(Q,P)$. Finally the triangle inequality for $b^\tau$ gives
$B_\tau(C)\le\mu(C)Z^\tau_a(C)$.
\end{proof}

A cut graph for $F$ uses pair capacities $k^\tau_{xy}$ and capacities
$u^{\tau,a}_x+\gamma\mu_xb^\tau_{ax}$ on edges $xt$, with $a$ inside
and $t$ outside. Replace $k^\tau$ by $\bar k^\tau$ of
Lemma~\ref{lem:prefix-core} and denote the result by $\bar F$.
Then $F\le\bar F\le2F$ on every set.

\begin{lemma}[Mass penalty scan]\label{lem:mass}
Suppose a cut procedure has factor $\alpha\ge1$ against every
narrow-balanced pinned comparator and returns a wide-balanced pinned set.
For a nonnegative cut objective $f$ and a bound $U>0$, a geometric scan
of penalties $\lambda\mu(C)$ produces, for every eligible $C^*$ with
$\mu(C^*)\le U$, a candidate satisfying
\[
 \mu(C)\le2\alpha U,\qquad f(C)\le3\alpha f(C^*).
\]
The scan has polylogarithmic length under the encoding assumptions.
\end{lemma}
\begin{proof}
Its comparison inequality is
$f(C_\lambda)+\lambda\mu(C_\lambda)
 \le\alpha(f(C^*)+\lambda U)$.
Let $\Delta$ be the smallest positive graph coefficient of $f$ and let
$f_{\max}$ be their sum. Start at $\lambda_0=\Delta/(4\alpha U)$
and double until the first returned set has mass at most $2\alpha U$.
Stop unsuccessfully after passing $2f_{\max}/U$; an eligible comparator
guarantees a passing set by then.
If the initial value passes, its cost is at most
$\alpha f(C^*)+\Delta/4$. When $f(C^*)>0$ this is at most
$3\alpha f(C^*)$, since $f(C^*)\ge\Delta$. When $f(C^*)=0$,
the output cost is less than $\Delta$ and therefore zero, and its mass
is at most $\alpha U$.
Otherwise the failed preceding value at $\lambda/2$ has mass greater
than $2\alpha U$. Its comparison inequality implies
$\lambda U<2f(C^*)$, which gives $f(C_\lambda)<3\alpha f(C^*)$.
If all coefficients vanish, use $\lambda=1$.
The ratio of endpoints of the scan has a polylogarithmic-bit logarithm.
Each penalty is represented by adding $\lambda\mu_x$ to capacity $xt$.
\end{proof}

Fix a global upper bound $\alpha$ for the chosen cut procedure.
Use Lemma~\ref{lem:fast-cut} for the fast implementation.
The polynomial-time combinatorial implementation instead uses
Lemma~\ref{lem:tree-dp}, amplified to factor $\alpha=O(\log n)$.
Its exact narrow balance also satisfies the stated wide balance.

For a state with $N>16$, enumerate all atoms of mass greater than one.
In addition take $C_1\log(n+1)$ independent uniform samples from its
unit-mass atoms, or enumerate those atoms if fewer exist. These are the
\emph{certificate anchors}. For each anchor $a$, each
$\lambda_a\le\tau\le H$, and each dyadic mass bound
$U\in\{1,2,4,\ldots,2^{\lceil\log_2n\rceil}\}$, do the following:
\begin{enumerate}[label=(\arabic*)]
\item Set $\gamma=2\alpha U$ and build the graph for
$f=\bar F_{\tau,a,\gamma}$.
\item Pin $a$ inside and $t$ and every atom with $\lambda_x>\tau$ outside.
Use the mass penalty scan of Lemma~\ref{lem:mass}, comparing to counts
in $[N/3,2N/3]$ and accepting counts only in $[N/8,7N/8]$.
\item Retain a returned set only when its mass is at most $2\alpha U$.
\end{enumerate}
Select the retained candidate of least $\bar F$ value.
If no candidate is valid, mark the trial failed.
After selection draw a fresh mass-proportional execution anchor and
use the recursion of Lemma~\ref{lem:l1-recursion}.
Small states use the enumerations after Lemma~\ref{lem:slow-separator},
with $\mathcal H$ in place of $\Gamma$.

\begin{lemma}[Few anchors and separator bound]\label{lem:fast-separator}
At depth $j$ there are at most $j$ non-unit-mass atoms.
There are $O(\log n)$ certificate anchors per state. Except with
inverse-polynomial probability, for every fixed current comparator $P$,
the selected candidate satisfies
\[
 \bar F(C)\le168\alpha^2D_\mu(Q,P),\qquad
 \E_A D_\mu(Q,T_AQ)\le336\alpha^2D_\mu(Q,P).
\]
\end{lemma}
\begin{proof}
The root has no non-unit atom. An inside child takes a subset; an outside
child introduces one new contracted atom. Thus the number increases
by at most one per depth. The wide balance gives child sizes at most
$7N/8+1\le15N/16$ for $N>16$, so depth is $O(\log n)$.

Choose the compatible $C^*,\tau^*$ of Lemma~\ref{lem:balanced}, write
$c=\mu(C^*)$ and $B=D_\mu(Q,P)$, and choose $c\le U<2c$.
For a mass-proportional anchor in $C^*$,
\[
 \E_a Z^{\tau^*}_a(C^*)=2B_{\tau^*}(C^*)/c.
\]
Together with Lemma~\ref{lem:unary} and $\gamma<4\alpha c$ this yields
$\E_a F_{\tau^*,a,\gamma}(C^*)\le14\alpha B$.
Call an anchor good when $F_{\tau^*,a,\gamma}(C^*)\le28\alpha B$.
Good anchors have mass at least $c/2$, and cut preservation implies
$\bar F(C^*)\le56\alpha B$ for each good anchor. This good-anchor
set is defined by exact coefficients and is independent of the graph
sampling used to approximate them. For $B=0$ every positive-mass
anchor is good and has zero score.
If a good anchor is non-unit it is enumerated. Otherwise there are at
least $c/2\ge N/6$ good unit atoms, and each uniform unit-atom sample
hits one with probability at least $1/6$.
The chosen logarithmic sample size suffices.
For its enumerated tuple Lemma~\ref{lem:mass} gives mass at most $\gamma$
and cost at most $168\alpha^2B$. Selection cannot increase the score.
The execution bound follows from
$\Gamma\le\mathcal H\le F\le\bar F$ and
Lemma~\ref{lem:certificate}.
\end{proof}

\begin{lemma}[Representing all recursive guide partitions]\label{lem:reps}
Assign to each atom $x$ an original representative $r_x\in V$.
Initially $r_x=x$. At a real layer $1\le i\le q$ its guide label is
\[
 \ell_i(x)=
 \begin{cases}
  \mathtt{all},&i\ge H,\\
  (\mathtt{private},x),&i<\lambda_x,\\
  (\mathtt{root},g_{\theta_{i-1}}(r_x)),&\lambda_x\le i<H.
 \end{cases}
\]
The tags make the three types distinct. The inside child inherits
representatives and upper index $\tau$. The contracted atom in the
outside child has representative $r_A$ and release $\tau$.
These labels give exactly the projected guide sequence.
One layer costs $\wt O(N)$ time and $O(N)$ additional space to report.
\end{lemma}
\begin{proof}
Induct on recursive operations. Inside restriction keeps fine relations
and becomes indiscrete at its upper index. Outside contraction is private
below $\tau$. At $\tau\le i<H$ the anchor is released, so its old label is
its root label, and the contracted atom joins exactly its old outside
anchor-block atoms. Unreleased outside atoms have private labels and remain
singletons. At $i\ge H$ everything is indiscrete. This is precisely
the projection definition. Two predecessor queries evaluate each root
label. Group equal labels by sorting. No ancestor partitions are stored.
\end{proof}

\begin{theorem}[The fast $L_1$ algorithm and the combinatorial one-pass $L_1$ algorithm]\label{thm:direct-l1}
The fast implementation has approximation $O(\log^5 n)$ in offline
time $\wt O(n^2)$ and in one-pass space $\wt O(n)$ with total update and
postprocessing time $\wt O(n^2)$ each.
The combinatorial implementation has approximation $O(\log^3 n)$
in one-pass space $\wt O(n)$ and polynomial time.
\end{theorem}
\begin{proof}
By Lemmas~\ref{lem:fast-separator} and~\ref{lem:l1-recursion}, expected
distance from the guide to the output is at most $O(\alpha^2\log n)$
times distance from the guide to a best feasible nested sequence.
To state the transfer to the input explicitly, fix an optimal
ultrametric $u^*$. Lemma~\ref{lem:grid} supplies a feasible grid sequence
$\widehat P$ with $D(Q,\widehat P)\le D(Q,P_{u^*})$. Hence
Lemma~\ref{lem:guide} and the triangle inequality give
\[
 \E L_1(d,R)
 \le \E D(G,Q)+\E D(Q,R)
 \le 7\OPT_1+O(\alpha^2\log n)\,8\OPT_1
 =O(\alpha^2\log n)\OPT_1,
\]
where $R$ is the output ultrametric, identified in $D(Q,R)$ with its
threshold partition sequence. Analyze an ideal execution that replaces a failed
separator by an existing valid one and then draws its fresh execution
anchor. There are polynomially many calls; each has a small conditional
failure probability. Coupling the implementation and ideal execution
until the first failure makes divergence inverse-polynomial.
This argument does not condition an anchor on future successful calls.
The same coupling treats record caps and graph compression failures.

Use $K=O(\log n)$ independent trials. Markov's inequality gives constant
success per ideal trial; the coupling preserves it up to inverse-polynomial
error. Select by exact cost offline, or by Lemma~\ref{lem:cost-sketch}
with $\varepsilon=1/8$ in streaming. The input sketch is independent
of the trials and is valid simultaneously for all candidates, so it can
be shared. This gives high-probability approximation.
If $\OPT_1=0$, the relevant nonnegative ideal expectations are zero,
so this argument includes that case.

For the fast time bound use the record grid, $q=\wt O(n)$.
At a state, accumulate the suffix arrays $b^\tau_{ax}$ and
$u^{\tau,a}_x$ for all $O(\log n)$ tried anchors by one sweep.
A forward sweep subtracts one layer at a time from these suffix arrays
and adds that layer's product cliques to the prefix graph of
Lemma~\ref{lem:prefix-core}. Each layer costs $\wt O(N)$.
At each birth, all mass guesses and penalty scans give only
polylogarithmically many cut calls on graphs with $\wt O(N)$ edges.
Thus state work is $\wt O(qN)$ and workspace is $\wt O(N)$.
At any depth, the potential $N-1$ is additive, since child counts are
$k,N-k+1$. The sum of nontrivial state counts at that depth is at most
$2(n-1)$. Hence $\sum_vN_v=O(n\log n)$, and total work is
$\wt O(qn)=\wt O(n^2)$.

Process states depth first, discarding working arrays before forming
children. The atom lists on the current path and deferred siblings use
$O(n)$ space because counts decrease geometrically apart from constant
small-state depth. Original record maps, the grid, current cut graphs,
suffix arrays, and output trees use $\wt O(n)$ words.
Before the pass choose the trial permutations, update their lists at
both endpoints of each record, and update $H_{\mathrm{in}},D_0,M$,
where $H_{\mathrm{in}}$ is the graph of Theorem~\ref{thm:prefix}.
Lemmas~\ref{lem:records} and~\ref{thm:prefix} give total update time
$\wt O(n^2)$ and the stated space. All subsequent guide queries use
stored representatives.

For the combinatorial implementation use the same procedure and arrays,
but the polynomial-time, linear-space tree cut procedure of
Lemma~\ref{lem:tree-dp}. Its $\alpha=O(\log n)$ gives the factor
$O(\log^3 n)$. Every graph passed to it has $\wt O(N)$ edges.
There are polynomially many calls, and its space is $\wt O(N)$.
No random computation is replayed. The explicit clique sampling in
Lemma~\ref{lem:prefix-core} uses arithmetic and sampling only.
\end{proof}


\section{A compact matrix for squared error}
\label{sec:l2-input}

Use the same permutation, functions, and records as in
Section~\ref{sec:guide}. Define a symmetric matrix with zero diagonal by
\begin{equation}\label{eq:bar-matrix}
 \bar d_{xy}=\int_0^M\ind[g_t(x)\ne g_t(y)]\,dt.
\end{equation}
Thus $0\le\bar d_{xy}\le M$. The matrix is not required to be an
ultrametric. For each vertex $r$ define
\[
 \beta_r=\min\{d_{rw}:\pi(w)<\pi(r)\},
\]
with the minimum of the empty set interpreted as $+\infty$.
Then $r$ is a center exactly for $t<\beta_r$. The value $\beta_r$ is
the key of the first record in the list for $r$ whose selected vertex
is not $r$, or $+\infty$ if that record does not exist. Thus it is
obtained from stored data.
On a record interval where $f_t(v)=r$, its label is $r$ before $\beta_r$
and $v$ afterwards. Splitting that interval at $\beta_r$ if necessary
therefore gives $O(\log n)$ intervals of constant $g_t(v)$ on every
successful capped execution. The distance $\bar d_{xy}$ is evaluated
by merging the two lists of intervals and summing the lengths on which
their labels differ. This costs $O(\log n)$ time and constant additional
workspace. Every entry is an integer in the arithmetic model of
Section~\ref{sec:definitions}. No dense matrix is needed.

Fix an ultrametric comparator $P$ independently of $\pi$, with entries
in $[0,M]$, and put $B=L_2(d,P)$. An \emph{input error interval} of a pair
$e$ is the interval between $d_e$ and $P_e$, with its endpoints removed
or added arbitrarily; it has length $|d_e-P_e|$.


Put $L=\lceil C_2\log(n+1)\rceil$, for a sufficiently large constant
$C_2$. For $v\in V$, let $F_v=\{f_t(v):0\le t<M\}$ and let $R_v$
be the set of minimum-ranked vertices of the comparator blocks containing
$v$, over $0\le t<M$. Define $\mathcal E$ to be the event that all input-prefix
record lists meet their caps and $|F_v|,|R_v|\le L$ for every $v$.

\begin{lemma}[Compact squared-error guide]\label{l2:lem:denoise}
The event $\mathcal E$ holds with high probability. On this event the stored records agree with the construction without record caps, and
\[
 \E_\pi\bigl[\ind[\mathcal E] L_2(\bar d,P)\bigr]
 \le O(\log^2 n)B.
\]
The event is used through its indicator, not by conditioning the permutation distribution.
\end{lemma}

\subsection{A witness for every disagreement}

At threshold $t$, the comparator block containing $v$ is $\{x:P_{vx}\le t\}$. Write $Q_t=Q(t)$.
An input pair $e$ is erroneous at $t$ when $t$ lies between $d_e$ and
$P_e$. Its erroneous-time interval has length $|d_e-P_e|$.
Every pairwise disagreement between $Q_t$ and the comparator partition
has such an input witness of one of four forms.

If $x,y$ are in a comparator block $C$ but have different guide labels,
let $r=\arg\min_{v\in C}\pi(v)$.
\begin{enumerate}[label=(\roman*),leftmargin=*]
\item If $r$ is a center, choose $v\in\{x,y\}$ whose label is not $r$.
If $d_{vr}>t$, witness $vr$: $P_{vr}\le t<d_{vr}$.
\item Still suppose $r$ is a center, but $d_{vr}\le t$. Choose the same endpoint $v$ with label different from $r$. Its selected neighbor $w=f_t(v)$ has
$\pi(w)<\pi(r)$, so $w\notin C$ and $d_{vw}\le t<P_{vw}$.
Witness $vw$.
\item If $r$ is not a center, let $w_r=f_{\beta_r}(r)$, its fixed
first lower-ranked selected neighbor. Then $w_r\notin C$,
$\pi(w_r)<\pi(r)$, and $d_{rw_r}=\beta_r\le t<P_{rw_r}$.
Witness $rw_r$. This one witness works for every later failure of $r$.
\end{enumerate}
If $x,y$ have the same guide label $r$ but are in different comparator
blocks, $r$ is a center and both are positively adjacent to it. At
least one endpoint $v$ is outside the comparator block of $r$.
This is the fourth form: witness $vr$, with $d_{vr}\le t<P_{vr}$.
For each unordered pair and each interval of the finite common input and comparator grid, choose the first applicable form and the lexicographically first possible witness within that form. Let $W_{xy}$ be the set of
witnesses selected for output pair $xy$ over all times.

\subsection{Few witnesses per output pair}

Comparator blocks containing $v$ form a fixed nested chain. Ordering
their newly added vertices by fixed ID reduces its minima to the same
left-to-right record calculation as above. Thus, simultaneously,
\[
 |F_v|\le L,\quad |R_v|\le L\qquad(v\in V)
\]
except with inverse-polynomial probability, by the exponential-moment argument of Lemma~\ref{lem:records}. This also proves the asserted probability of $\mathcal E$.
The four witness forms lie among
\[
 \{vr:r\in R_v,\ v\in\{x,y\}\},\quad
 \{vw:w\in F_v,\ v\in\{x,y\}\},\quad
 \{rw_r:r\in R_x\cup R_y,\ \beta_r<+\infty\}.
\]
Consequently $|W_{xy}|\le O(L)$ on $\mathcal E$.
Integration and Cauchy--Schwarz give
\begin{equation}\label{l2:eq:witness-square}
 |\bar d_{xy}-P_{xy}|^2
 \le |W_{xy}|\sum_{e\in W_{xy}}|d_e-P_e|^2.
\end{equation}

\subsection{Few output pairs in expectation per input witness}

Write $A_e=|\{xy:e\in W_{xy}\}|$. We show
$\E A_e=O(\log n)$, without conditioning on $\mathcal E$.
All sizes below are maximized over the indicated nested block chains.

For form (i), orient $e=vr$. All charged output pairs are $vy$ with
$y$ in a comparator block containing $r$ whose minimum is $r$.
Let $T_r$ be the largest size of such a block. If $T_r\ge k$, $r$ is
minimum in the first block of its chain of size at least $k$; hence
$\Prb[T_r\ge k]\le1/k$, and $\E T_r\le1+\log n$.

For form (ii), orient $e=vw$. Charged pairs are $vy$ inside a block
$C$ containing $v$, excluding $w$, with $w$ preceding every vertex of
$C$. For the largest such size $T_{v,w}$,
\[
 \Prb[T_{v,w}\ge k]\le1/(k+1),\qquad
 \E T_{v,w}=O(\log n).
\]
Indeed, at the first eligible fixed block of size at least $k$, $w$
must be first among that block and $w$.

For form (iii), orient $e=rw$. Charged output pairs lie inside a block
$C$ containing $r$, excluding $w$, where $r$ is minimum in $C$ and
$w$ precedes $r$. If $T$ is the largest such block size, the first
eligible fixed block of size $s\ge k$ requires $w$ first and $r$ second
among its $s+1$ vertices. Therefore
\[
 \Prb[T\ge k]\le\frac1{k(k+1)},\qquad
 \E T^2=\sum_{k=1}^n(2k-1)\Prb[T\ge k]=O(\log n).
\]
There are at most $T^2/2$ charged unordered output pairs. An empty eligible chain is assigned size zero in each of these maximum-size definitions.

For the fourth form, orient $e=vr$. Every charged pair is $vy$ inside
the closed positive neighborhood of center $r$. These neighborhoods
grow in a fixed order as $t$ increases, and $r$ must be their minimum.
Define $K_r$ as the maximum size of this neighborhood while $r$ is
a center, and put $K_r=0$ if there is no such threshold. Then
$\Prb[K_r\ge k]\le1/k$, so $\E K_r=O(\log n)$.
Summing the four bounds and both orientations proves $\E A_e=O(\log n)$.

\begin{proof}[Conclusion of Lemma~\ref{l2:lem:denoise}]
Sum \eqref{l2:eq:witness-square} over output pairs on $\mathcal E$ and
then take expectation. Since all terms are nonnegative,
\[
 \E[\ind[\mathcal E]L_2(\bar d,P)]
 \le O(L)\sum_e |d_e-P_e|^2\E A_e
 =O(\log^2 n)B.
\]
The comparator is fixed independently of $\pi$. Random grid endpoints
are not conditioned on, and the event's indicator is never moved
inside a conditional permutation distribution.
\end{proof}

\section{The squared-error recursion on an implicit matrix}

This section applies to any symmetric matrix $\mathsf D^0$ with zero diagonal and entries in $[0,M]$, evaluable per pair in polylogarithmic time, including $\mathsf D^0=\bar d$. The term pair query means evaluation of one specified matrix entry.
At a state, let $X$ be its $N$ atoms, with integer masses $\mu_x\ge1$,
real release heights $0\le\lambda_x\le H$, upper height $0\le H\le M$, and original representatives $r_x\in V$. Release heights in this section are real heights, not the layer indices of Section~\ref{sec:projection}.
Initially $\mu_x=1$, $\lambda_x=0$, $H=M$, and $r_x=x$.
The state input is
\begin{equation}\label{l2:eq:state-input}
 \mathsf D_{xy}=\min\{H,\max\{\lambda_x,\lambda_y,\mathsf D^0_{r_xr_y}\}\}
 \quad(x\ne y),\qquad \mathsf D_{xx}=0.
\end{equation}
A feasible comparator is an ultrametric $P$ with
$\max\{\lambda_x,\lambda_y\}\le P_{xy}\le H$ for $x\ne y$.
For state matrices use the weighted squared distance
\[
 B_\mu(\mathsf D,P)=\sum_{xy\in\binom X2}\mu_x\mu_y(\mathsf D_{xy}-P_{xy})^2.
\]
Total mass is at most $n$. Atom-count balance is always a condition
on $|C|$, independently of $\mu(C)$.

\subsection{Clamped projection and representatives}

For $\varnothing\ne C\subsetneq X$, $c=\mu(C)$,
$\max_{x\in C}\lambda_x\le\tau\le H$, and $A\in C$, define
\[
 (T_{C,\tau,A}\mathsf D)_{xy}=
 \begin{cases}
 \min(\mathsf D_{xy},\tau),&x,y\in C,\\
 \max(\mathsf D_{Ay},\tau),&x\in C,\ y\notin C,\\
 \mathsf D_{xy},&x,y\notin C,
 \end{cases}
\]
with symmetry in the cross case and zero diagonal. We call $\tau$ the
birth height of $C$. As in Section~\ref{sec:projection}, a separator
is a pair $(C,\tau)$ with $2\le |C|\le N-1$ satisfying the displayed
release and upper-height inequalities. In this section $T$ abbreviates $T_{C,\tau,A}$ only when the set, height, and execution anchor are fixed.
The inside child has the atoms of $C$, unchanged masses, releases,
representatives, and upper height $\tau$. The outside child replaces
$C$ by $z$ of mass $c$, release $\tau$, representative $r_A$, and
inherits upper height $H$. Formula \eqref{l2:eq:state-input} is preserved:
the inside update clips at $\tau$, and the outside cross update is
$\max\{\mathsf D_{Ay},\tau\}$, with $\lambda_A\le\tau$.

\begin{lemma}[Squared nonexpansion]\label{l2:lem:nonexpand}
For arbitrary state matrices $\mathsf D,P$, and a fresh anchor with
$\Prb[A=x]=\mu_x/c$,
\[
 \E_A B_\mu(T\mathsf D,TP)\le B_\mu(\mathsf D,P).
\]
If $P$ is feasible, $TP$ is a feasible ultrametric for the two children
and their substitution. Their weighted errors sum to $B_\mu(T\mathsf D,TP)$.
\end{lemma}
\begin{proof}
For real $z,z^\prime$, both $|\min(z,\tau)-\min(z^\prime,\tau)|$ and $|\max(z,\tau)-\max(z^\prime,\tau)|$ are at most $|z-z^\prime|$. Internal errors decrease.
For each outside $y$, the expected cross error is at most
\[
 \E_A[c\mu_y(\mathsf D_{Ay}-P_{Ay})^2]
 =\sum_{x\in C}\mu_x\mu_y(\mathsf D_{xy}-P_{xy})^2.
\]
Outside errors are unchanged. For ultrametricity, internal distances
are at most $\tau$, and each outside vertex has a common distance to
$C$, at least $\tau$. For outside $y,w$, the ultrametric inequalities
for the triple $A,y,w$ remain valid after flooring both anchor
distances at $\tau$. Releases and upper bounds follow from the stated
birth constraint. Internal, cross, and outside pair scopes are disjoint,
which proves additivity of child errors.
\end{proof}

Call $C$ compatible with $P$ at $\tau$ when internal distances are at
most $\tau$, cross distances are at least $\tau$, and, for every outside
$y$, $P_{xy}$ is constant over $x\in C$. Then $TP=P$ for every anchor, using the projection defined in this section.
Every feasible ultrametric with $N>16$ has such a set with
$N/3\le |C|\le2N/3$. Descend through a child of size greater than
$2N/3$ until none exists; choose a child in $[N/3,2N/3]$, or accumulate
children smaller than $N/3$ until their union first reaches $N/3$.
Use that node's height as the birth for a union, or the chosen child's
height. Releases are at most this height. For $3\le N\le16$, a proper
compatible set with $2\le |C|\le N-1$ exists by the same hierarchy argument.

\subsection{A variance certificate represented by a cut}

Put $h^\tau_{xw}=\max(\mathsf D_{xw},\tau)$, including $h^\tau_{xx}=\tau$,
and define
\begin{align*}
 b^{\tau,2}_{xa}&=\sum_{w\in X}\mu_w(h^\tau_{xw}-h^\tau_{aw})^2,\\
 k^{\tau,2}_{xy}&=\mu_x\mu_y(\tau-\mathsf D_{xy})_+^2,\\
 u^{\tau,a,2}_x&=\mu_x b^{\tau,2}_{xa},\qquad
 Z^{\tau,2}_a(C)=\sum_{x\in C}\mu_x(\mathsf D_{ax}-\tau)_+^2,\\
 \mathcal J_{\tau,a}(C)&=k^{\tau,2}(\cut C)+u^{\tau,a,2}(C)+\mu(C)Z^{\tau,2}_a(C).
\end{align*}
For $\gamma\ge\mu(C)$, the
graph-cut objective
\begin{equation}\label{l2:eq:F}
 F^{(2)}_{\tau,a,\gamma}(C)
 =k^{\tau,2}(\cut C)+u^{\tau,a,2}(C)+\gamma Z^{\tau,2}_a(C)
\end{equation}
dominates $\mathcal J$. Pair edges have capacities $k^{\tau,2}_{xy}$;
an edge from $x$ to a forced-out dummy sink has capacity
$u^{\tau,a,2}_x+\gamma\mu_x(\mathsf D_{ax}-\tau)_+^2$.

\begin{lemma}[Certificate dominates expected squared edit]\label{l2:lem:edit}
For every certificate anchor $a\in C$ and fresh execution anchor
$\Prb[A=x]=\mu_x/c$,
\[
 \E_A B_\mu(\mathsf D,T_{C,\tau,A}\mathsf D)\le6\mathcal J_{\tau,a}(C).
\]
\end{lemma}
\begin{proof}
For internal pairs,
$(\mathsf D_{xy}-\tau)_+\le|h^\tau_{xy}-h^\tau_{ay}|+(\mathsf D_{ay}-\tau)_+$.
Write $k=k^{\tau,2}(\cut C)$, $u=u^{\tau,a,2}(C)$, and
$Z=Z^{\tau,2}_a(C)$ for this proof. Squaring and summing over unordered
internal pairs bounds the internal edit by $2u+2cZ$.
For cross pairs, first floor $\mathsf D_{xy}$ at $\tau$, giving at most
$2k^{\tau,2}(\cut C)$ plus twice the row-difference contribution.
For each outside $y$, put
$m_y=c^{-1}\sum_{x\in C}\mu_xh^\tau_{xy}$.
Expanding squares gives
$\E_A\sum_{x\in C}\mu_x(h^\tau_{xy}-h^\tau_{Ay})^2
=2\sum_{x\in C}\mu_x(h^\tau_{xy}-m_y)^2$.
The choice $m_y$ minimizes the sum of squared deviations, so
\[
 \E_A\sum_{x\in C}\mu_x(h^\tau_{xy}-h^\tau_{Ay})^2
 \le2\sum_{x\in C}\mu_x(h^\tau_{xy}-h^\tau_{ay})^2.
\]
Thus cross edit is at most $2k+4u$. The total is at most
$2k+6u+2cZ\le6\mathcal J$.
\end{proof}

\section{Only linearly many birth candidates per anchor}

For a certificate anchor $a$, use
\[
 \mathcal T_a=\{0,H\}\cup\{\lambda_x:x\in X\}
                 \cup\{\mathsf D_{ax}:x\in X\}.
\]
Its size is $O(N)$. Unlike a grid of all pair distances, this keeps
the number of cut calls nearly linear per state.

\begin{lemma}[Rounding to an anchor row]\label{l2:lem:round}
Let $C$ be compatible with $P$ at $\tau_0$, $c=\mu(C)$, and
$B=B_\mu(\mathsf D,P)$. For each $a\in C$, let $\tau_a$ be the nearest point
of $\mathcal T_a$ to $\tau_0$, choosing the smaller point in a tie. Then it is a valid birth for $C$ and
\[
 \E_{a\sim\mu|_C}\mathcal J_{\tau_a,a}(C)\le54B.
\]
\end{lemma}
\begin{proof}
If $\tau_0$ is not a grid point, let $l<\tau_0<r$ be its adjacent
grid points. No $\mathsf D_{aw}$ lies in $(l,r)$. All included releases lie
at most at $l$, since they are grid points, and $r\le H$.
Let $e_{aw}=\mathsf D_{aw}-P_{aw}$ and let $P'$ be the projection of $P$ at
$\tau_a$, which is independent of execution anchor because the old
cross distances are constant on $C$.

If rounding down to $l$, an affected internal pair has
$l<P_{xy}\le\tau_0$. Ultrametricity gives an endpoint $z\in\{x,y\}$
with $P_{az}\ge P_{xy}$. If $\mathsf D_{az}\le l$, its error is at least
$P_{xy}-l$. If $\mathsf D_{az}\ge r$, its error is at least
$r-\tau_0\ge\tau_0-l\ge P_{xy}-l$.
If rounding up to $r$, an affected cross pair has
$\tau_0\le P_{xy}<r$ and $P_{ay}=P_{xy}$. If $\mathsf D_{ay}\ge r$, its
error is at least $r-P_{xy}$. If $\mathsf D_{ay}\le l$, its error is at least
$\tau_0-l\ge r-\tau_0\ge r-P_{xy}$.
These arguments also apply to the affected anchor-row entries. Hence
\begin{equation}\label{l2:eq:round-cost}
 B_\mu(P,P')\le c\sum_w\mu_w e_{aw}^2,
 \qquad |P'_{aw}-P_{aw}|\le|e_{aw}|.
\end{equation}
When $\tau_0$ is itself a grid point these bounds are immediate.
Taking the mass-proportional average of the first bound gives at most
$2B$, and therefore $\E B_\mu(\mathsf D,P')\le6B$.

Since $P'$ is compatible at $\tau_a$, its cross distances are at least
$\tau_a$, so $k^{\tau_a,2}(\cut C)\le B_\mu(\mathsf D,P')$.
Its floored rows on $C$ are identical. Triangle inequality squared
therefore gives
\[
 u^{\tau_a,a,2}(C)\le4B_\mu(\mathsf D,P')
       +8c\sum_w\mu_w e_{aw}^2.
\]
Also $cZ^{\tau_a,2}_a(C)\le4c\sum_{x\in C}\mu_xe_{ax}^2$.
The expected bounds on the three terms are $6B$, $40B$, and $8B$.
Their sum proves the claim.
\end{proof}

\section{Sparse cut instances and the separator algorithm}

\subsection{The cut and mass procedures}
Use Lemma~\ref{lem:fast-cut} with its factor $\alpha=O(\log^2 n)$,
comparing to counts in $[N/3,2N/3]$ and accepting counts in $[N/8,7N/8]$.
An atom with $\lambda_x>\tau$ is forced outside by contracting it into the
additional sink and retaining its count. All returned constraints are checked.
For an intended mass bound $U$, Lemma~\ref{lem:mass} applied to the cut
objective $f$ gives
\begin{equation}\label{l2:eq:mass-penalty}
 \mu(C)\le2\alpha U,\qquad f(C)\le3\alpha f(C^*)
\end{equation}
for every eligible $C^*$ of mass at most $U$.

\subsection{Sparsifying the squared-error cut graph}

Generate the $N(N-1)/2$ state pairs through \eqref{l2:eq:state-input},
with keys $\mathsf D_{xy}$ and capacities $\mu_x\mu_y$.
Construct a $(1\pm1/4)$ all-threshold cut sketch $K$ with
$\wt O(N)$ retained edges. Each retained edge keeps its key and its
new integer capacity. Scale by $4/3$ for domination.
For every $\tau$, give a retained edge of key $z$ and scaled capacity
$a$ the capacity $a(\tau-z)_+^2$. This graph satisfies simultaneously
\begin{equation}\label{l2:eq:hinge-sketch}
 k^{\tau,2}(\cut C)\le\bar k^{\tau,2}(\cut C)\le2k^{\tau,2}(\cut C).
\end{equation}
Indeed,
$(\tau-z)_+^2=\int_0^\tau2(\tau-t)\ind[z\le t]\,dt$.
Integrate the nonnegative prefix-cut inequalities against this kernel.

Use Theorem~\ref{thm:prefix} with vertices $X$ and global failure
parameter $n^{-a}$ for a sufficiently large fixed $a$. The number of
sampled forests uses $\log n$, so it also covers small states. Its
resource bound is $\widetilde O(N)$ space and $\widetilde O(1)$ update
time per generated unit-capacity pair.

For weighted pairs, write $\mu_x\mu_y$ in binary and run the unweighted
routine on each bit layer, scaling its output by that bit's power of
two. Summing bit-layer inequalities preserves the same relative error.
There are $O(\log n)$ layers. Generate their pairs again from the
implicit matrix, instead of storing a dense input. This gives
$\wt O(N^2)$ time and $\wt O(N)$ space for the state sketch, with
inverse-polynomial failure probability. Generated pair capacities have $O(\log n)$ bits; sampling exponents also have $O(\log n)$ bits, so the resulting capacities have polylogarithmic encodings.

\subsection{Approximating row norms by random signs}

The random-sign Johnson--Lindenstrauss theorem says that a map
$z\mapsto r^{-1/2}(\sum_i\sigma_{ij}z_i)_{j=1}^r$, with independent
uniform signs $\sigma_{ij}\in\{-1,1\}$, preserves all squared Euclidean
norms in a fixed finite set of $n^{O(1)}$ vectors within $1\pm1/4$,
with prescribed inverse-polynomial failure, for $r=C\log n$.
Here $\|z\|_2^2=\sum_i z_i^2$. We use Theorem 1.1 of~\cite{Ach03}.
This is an analysis of explicit sign sampling, not a linear-system call.

Choose $r=O(\log n)$ fresh random-sign projection coordinates for the
state, enough for its prescribed failure probability. Conceptually
expand column $w$ into $\mu_w$ identical coordinates, and set
\[
 s_{w,j}=\sum_{i=1}^{\mu_w}\sigma_{w,i,j},
 \qquad
 V_{\tau,x,j}=\sum_{w\in X}s_{w,j}\max(\mathsf D_{xw},\tau),
 \qquad \sigma_{w,i,j}\in\{-1,1\}.
\]
Signs are independent. Only the integer aggregates $s_{w,j}$ are
stored. The random-sign Johnson--Lindenstrauss theorem \cite[Theorem 1.1]{Ach03} gives,
simultaneously for every tried $(\tau,a,x)$,
\[
 (3/4)b^{\tau,2}_{xa}\le
 \frac1r\sum_{j=1}^r(V_{\tau,x,j}-V_{\tau,a,j})^2
 \le(5/4)b^{\tau,2}_{xa}.
\]
Define
\[
 \bar b^{\tau,2}_{xa}=\frac{4}{3r}\sum_{j=1}^r
       (V_{\tau,x,j}-V_{\tau,a,j})^2,\qquad
 \bar u^{\tau,a,2}_x=\mu_x\bar b^{\tau,2}_{xa}.
\]
Thus $b^{\tau,2}_{xa}\le\bar b^{\tau,2}_{xa}\le2b^{\tau,2}_{xa}$.
Choose certificate anchors and their
grids before these fresh signs; there are only polynomially many fixed
vectors to preserve conditionally on the state input. The factor $r^{-1/2}$ is used only in the theorem: the implementation computes the displayed sum of squares divided by $r$, using rational arithmetic.

Define $\bar F^{(2)}$ by \eqref{l2:eq:F} using $\bar k^{\tau,2}$, $\bar u^{\tau,a,2}$, and exact $Z^{\tau,2}_a$.
For $\gamma\ge\mu(C)$ it dominates $\mathcal J_{\tau,a}(C)$. If
$\mu(C^*)\le U<2\mu(C^*)$ and $\gamma=2\alpha U$, then
\begin{equation}\label{l2:eq:surrogate-factor}
 \bar F^{(2)}_{\tau,a,\gamma}(C^*)\le4\alpha\mathcal J_{\tau,a}(C^*).
\end{equation}

\subsection{Concrete separator procedure and its bound}

For $N>16$, enumerate all atoms of mass greater than one, and, when
unit-mass atoms exist, sample $C_c\log n$ of them uniformly and
independently. They are
certificate anchors. For each anchor $a$, each $\tau\in\mathcal T_a$
with $\lambda_a\le\tau$, and each dyadic $U$ up to $2n$:
\begin{enumerate}[leftmargin=*]
\item Put $\gamma=2\alpha U$ and construct its sparse $\bar F^{(2)}$ graph.
\item Force $a$ in, the dummy sink out, and all $\lambda_x>\tau$ out.
Apply the geometric mass-penalty scan to this graph.
\item Retain a valid result only if $\mu(C)\le\gamma$.
Record its $\bar F^{(2)}$ value. After all calls, choose the cheapest
retained candidate, including its birth $\tau$.
\end{enumerate}
Then sample an independent execution anchor $A$ mass-proportionally
from $C$, and recurse on the two projected children.
For $3\le N\le16$, enumerate all proper $C$ with $2\le|C|\le N-1$,
all $a\in C$, and their valid anchor-grid births
($\max_{x\in C}\lambda_x\le\tau\le H$), selecting minimum exact
$\mathcal J$ before the independent execution draw. This is constant
work apart from coefficient arithmetic. For $N=2$, return its sole
input distance $\mathsf D_{12}$; for $N=1$, return its one leaf.

\begin{lemma}[Separator guarantee]\label{l2:lem:separator}
For every feasible comparator $P$, the chosen separator has
\[
 \E_A B_\mu(\mathsf D,T\mathsf D)\le O(\alpha^2)B_\mu(\mathsf D,P)
\]
except with the stated inverse-polynomial subroutine probability.
The execution draw remains fresh and independent of separator selection.
\end{lemma}
\begin{proof}
For a narrow-balanced compatible $C^*$, Lemma~\ref{l2:lem:round} gives
mass at least $\mu(C^*)/2$ of anchors $a$ with
$\mathcal J_{\tau_a,a}(C^*)\le108B$.
If one is a non-unit atom, it is enumerated. Otherwise at least
$\mu(C^*)/2\ge N/6$ are unit atoms, so the samples hit one with
inverse-polynomially high probability. A matching dyadic $U$ exists.
Equations \eqref{l2:eq:mass-penalty} and \eqref{l2:eq:surrogate-factor} give
a returned $\bar F^{(2)}$ value at most $12\alpha^2\mathcal J_{\tau_a,a}(C^*)$.
Its mass is at most $\gamma$, so its score dominates $\mathcal J$.
The minimum over candidates is no larger. Lemma~\ref{l2:lem:edit} now
bounds the independently executed squared edit.
For small states, an enumerated compatible set and an anchor attaining
the expectation bound of Lemma~\ref{l2:lem:round} give the same conclusion.
The number of non-unit atoms at depth $d$ is at most $d$: initially
none exist; an outside child adds only one, and an inside child retains
only a subset. Thus only $O(\log n)$ anchors are tried per state.
\end{proof}

\section{Approximation of the recursive core and output selection}

For $N>16$, child counts are $|C|$ and $N-|C|+1$. They are at most
$15N/16$, since $N/8\le|C|\le7N/8$. Small states have only a constant
number of further levels. Let $h=O(\log n)$ be total depth.

\begin{theorem}[Squared-error core]\label{l2:thm:core}
For any initial matrix $\mathsf D^0$ and ultrametric comparator $P$, the ideal
recursion returns a feasible ultrametric $U$ with
\[
 \E L_2(\mathsf D^0,U)\le O(\alpha^2\log^2 n)L_2(\mathsf D^0,P).
\]
The implemented recursion agrees with it except with inverse-polynomial
probability. All selected births and output heights are integers when
$\mathsf D^0$ and the initial bounds are integers.
\end{theorem}
\begin{proof}
Project the comparator with the same independently drawn execution
anchors as the state input. Lemma~\ref{l2:lem:nonexpand} shows that expected
total comparator error over states at any depth is at most its root
value. Lemma~\ref{l2:lem:separator} therefore bounds expected total squared
edit at one depth by $O(\alpha^2)L_2(\mathsf D^0,P)$.

Lift each state to its original pair scope: inside-child pairs, cross
pairs represented by the contracted atom, and outside pairs are
disjoint. Mass products are exactly the counts of represented pairs.
One original pair follows one child scope at each depth. If its successive
edits are $\Delta_1,\ldots,\Delta_h$, then
$(\sum_i\Delta_i)^2\le h\sum_i\Delta_i^2$.
Summing over original pairs and depths gives the $O(\alpha^2h^2)$ bound.
At termination, substitution of inside trees into contracted leaves
is feasible because internal heights are at most each birth and
outside distances are at least that birth. Its final pair values are
precisely the terminal edited matrix values.

For the failure coupling, run an ideal process that supplies a separator
meeting the proved bound whenever a certificate search, sketch, random-sign map,
or cut call fails, and then uses a fresh execution draw. Each repair
exists by the preceding lemmas. Up to the first such event, the ideal
and actual processes agree. There are polynomially many calls, each
with small conditional failure probability, so a union bound couples
them with inverse-polynomial failure. No argument conditions an
execution anchor on future subroutine success. Small-state substitutions
and integer anchor grids preserve integer heights.
\end{proof}

\begin{theorem}[One trial and amplification]\label{l2:thm:approx}
With constant probability greater than $3/4$, one complete independent
trial on $\bar d$ returns $U$ with
\[
 L_2(d,U)\le O(\alpha^2\log^4 n)\OPT_2.
\]
$K=O(\log n)$ independent trials and the selection rule below amplify
this to any desired inverse-polynomial failure probability, without
access to the original input after its scan.
\end{theorem}
\begin{proof}
Fix an optimal original comparator $P$. Conditional on the permutation
and its implicit $\bar d$, Theorem~\ref{l2:thm:core} applies with this same
$P$. Its internal randomness is fresh. Lemma~\ref{l2:lem:denoise} gives
\[
 \E[\ind[\mathcal E]L_2(\bar d,U)]
 \le O(\alpha^2\log^2 n)
        \E[\ind[\mathcal E]L_2(\bar d,P)]
 \le O(\alpha^2\log^4 n)\OPT_2.
\]
Use squared triangle inequality for $d,\bar d,U$ and for $d,P,\bar d$.
It bounds the corresponding truncated expectation of $L_2(d,U)$ by
the same order. Markov's inequality with a sufficiently large absolute
constant, together with the rare complement of $\mathcal E$ and the
core coupling failure, gives success probability greater than $3/4$.
For zero optimum, these truncated nonnegative expectations are zero,
so the statement still holds.

Choose $K$ odd. For outputs $U_1,\ldots,U_K$, evaluate all squared
distances $\Delta_{ij}=\sum_{x<y}(U_{i,xy}-U_{j,xy})^2$.
Return the output minimizing the median of its $K$ distances.
With high probability a strict majority lie within norm distance
$R=\sqrt{C\alpha^2\log^4 n\,\OPT_2}$ of $d$.
Every good output has median norm distance at most $2R$, because the norm distance between two good outputs is at most $2R$. The median of squared distances is minimized by exactly the same output as the median of norm distances, since square root is increasing. The selected
output has at least $(K+1)/2$ outputs within $2R$, one of which must
be good. Its distance from $d$ is at most $3R$, a factor $9$ in squared
error. Chernoff's bound supplies the desired high probability.
No input-cost approximation or second pass is used.
\end{proof}

\section{Time and space: one algorithm with different batch sizes}

Across a run, retain its compact label trajectories and its output tree.
Each original matrix query costs $O(\log n)$ on the successful capped
execution. Equation \eqref{l2:eq:state-input} then takes one such query.

At a state, choose its certificate anchors and form the union of their
grids. It has $q=\wt O(N)$ heights. Generate the sparse cut graph
once in $\wt O(N^2)$ time and $\wt O(N)$ space. Keep only the best
separator seen; each graph has $\wt O(N)$ edges, and there are
$\wt O(N)$ cut calls. Their total time is $\wt O(N^2)$.

\subsection{Computing threshold features in batches}

Let $\mathcal W\ge n$ be a workspace parameter, with logarithmic
factors absorbed in the space bound. Batch the sorted trial heights
into groups of
\[
 b=\min\{q,\max(1,\lfloor\mathcal W/N\rfloor)\}
\]
heights. One batch's $N\times b\times r$ feature table fits in
$\wt O(\mathcal W)$ space. It can be computed in one generated-pair
scan, using constant and slope difference arrays.

For a row $x$ and column $w$, let $z=\mathsf D_{xw}$. Its contribution to
$V_{\tau,x,j}$ is $s_{w,j}z$ for $\tau<z$ and $s_{w,j}\tau$ for
$\tau\ge z$. Binary-search the first batch height at least $z$.
Add its constant contribution before that index and its slope
contribution from that index onward, in $O(r\log b)$ work.
Prefix the difference arrays to obtain all batch features in
$O(Nbr)$ work. Include the diagonal $z=0$ explicitly. Queries within
the batch are then processed with the same fixed sketches and signs.
Discard the batch table and proceed to the next batch.

Thus one state needs
\[
 \wt O\left(N^2+\frac{N^4}{\mathcal W}+W_X\right)
 \quad\text{time},\qquad
 \wt O(\mathcal W+N)\quad\text{space},
\]
where $W_X=\mu(X)\le n$. The last work term samples the aggregated
signs; literal sampling of $W_Xr$ signs suffices and stores only their
atom sums. No dense pair list or feature table is needed in small space.
With large workspace, all $q$ heights form one batch.

\subsection{Summing the recursion resources}

The potential $N-1$ is additive across children. At any depth, total
count over nontrivial states is at most $2(n-1)$, while maximum count
decreases geometrically apart from constant-depth small states.
Lemma~\ref{lem:recursion-size} consequently gives
\[
 \sum_v N_v^2=O(n^2),\qquad
 \sum_v N_v^4=O(n^4).
\]
There are $O(n)$ states and each has mass at most $n$, so
$\sum_v W_{X_v}=O(n^2)$. The total core time is
\begin{equation}\label{l2:eq:tradeoff}
 \boxed{\wt O(n^2+n^4/\mathcal W)},
 \qquad\text{working space }\wt O(\mathcal W+n).
\end{equation}
After choosing a separator, discard parent working arrays and process
children depth first. Stored atom lists on the active path and deferred
siblings have $O(n)$ total size because their counts decrease
geometrically. Representatives, masses, and releases are the only
state records retained.

For amplification, all $K$ compact inputs and output trees occupy
$\wt O(n)$ words. A lowest-common-ancestor data structure for each hierarchy gives a pair
distance in constant time; alternatively, ancestor searches in logarithmic time preserve the displayed bounds. Scan original vertex pairs, query the $K$
trees, and accumulate their $K^2$ mutual squared distances. This takes
$\wt O(n^2)$ time and $\wt O(n)$ space. It does not query the original
input distances. Processing different trials sequentially preserves
\eqref{l2:eq:tradeoff}.

\section{Offline and one-pass squared-error guarantees}

\begin{theorem}[Offline and one-pass bounds]\label{l2:thm:models}
Under the arithmetic convention of Section~\ref{sec:definitions}, using Lemma~\ref{lem:fast-cut}, the offline
implementation returns the approximation of Theorem~\ref{l2:thm:approx}
in $\wt O(n^2)$ time and $\wt O(n^2)$ space.
Its one-pass implementation stores
$\wt O(n)$ words, takes $\wt O(n^2)$ total stream-update time, and
$\wt O(n^3)$ postprocessing time, with the same approximation.
\end{theorem}
\begin{proof}
Before reading input, choose the $K$ independent permutations. Maintain
their record maps and the global maximum $M$ in the same input scan.
Afterward, derive their label trajectories. All generated state pairs
come from these trajectories and representatives. The original stream
is never reread. Offline use $\mathcal W=n^2$; streaming use
$\mathcal W=n$ in \eqref{l2:eq:tradeoff}. All separator decisions, random
signs, execution anchors, and output selection are identical; a batch
partition changes only how their feature values are computed.
\end{proof}



\begin{proof}[Proof of Theorem~\ref{thm:four}]
The combinatorial offline assertion is Theorem~\ref{thm:slow}.
Both direct-cut $L_1$ assertions are Theorem~\ref{thm:direct-l1};
the combinatorial one-pass assertion is its polynomial-time instantiation.
For squared error, Theorems~\ref{l2:thm:approx} and~\ref{l2:thm:models}
give factor $O(\alpha(n)^2\log^4 n)$, which is $O(\log^8 n)$ for
$\alpha(n)=O(\log^2 n)$. Taking square roots gives $O(\log^4 n)$
for the norm error. All displayed space bounds include postprocessing.
All displayed probabilities include the caps, cut failures, sparsifier
failures, sign-projection failures, and amplification failures.
\end{proof}

\begin{thebibliography}{10}
\small
\bibitem{Ach03}
Dimitris Achlioptas.
Database-friendly random projections: Johnson--Lindenstrauss with binary coins.
\emph{Journal of Computer and System Sciences} 66(4), 671--687, 2003.
Theorem 1.1.
\url{https://doi.org/10.1016/S0022-0000(03)00025-4}

\bibitem{ADK25}
Daniel Agassy, Dani Dorfman, and Haim Kaplan.
Expander decomposition for non-uniform vertex measures.
arXiv:2510.23913, version 2, 2025.
Theorem 3.2 and the weighted-graph conventions of Section 2.
\url{https://arxiv.org/abs/2510.23913}

\bibitem{BK15}
Andr\'as A. Bencz\'ur and David R. Karger.
Randomized approximation schemes for cuts and flows in capacitated graphs.
\emph{SIAM Journal on Computing} 44(2), 290--319, 2015.
Theorems 5.2, 6.2, 6.4--6.5, 8.3, and footnote 2 following Theorem 8.3.
\url{https://doi.org/10.1137/070705970}

\bibitem{FRT04}
Jittat Fakcharoenphol, Satish Rao, and Kunal Talwar.
A tight bound on approximating arbitrary metrics by tree metrics.
\emph{Journal of Computer and System Sciences} 69(3), 485--497, 2004.
\url{https://doi.org/10.1016/j.jcss.2004.04.011}

\bibitem{LNPS23}
Jason Li, Danupon Nanongkai, Debmalya Panigrahi, and Thatchaphol Saranurak.
Near-linear time approximations for cut problems via fair cuts.
In \emph{Proceedings of SODA 2023}, 240--275.
Theorem 1.3.
\url{https://arxiv.org/abs/2203.00751}

\bibitem{Racke08}
Harald R\"acke.
Optimal hierarchical decompositions for congestion minimization in networks.
In \emph{Proceedings of STOC 2008}, 255--264.
Theorems 2 and 4, Lemmas 7--9.
\url{https://doi.org/10.1145/1374376.1374415}

\bibitem{SS11}
Daniel A. Spielman and Nikhil Srivastava.
Graph sparsification by effective resistances.
\emph{SIAM Journal on Computing} 40(6), 1913--1926, 2011.
Theorem 1.
\url{https://doi.org/10.1137/080734029}

\bibitem{Tropp12}
Joel A. Tropp.
User-friendly tail bounds for sums of random matrices.
\emph{Foundations of Computational Mathematics} 12, 389--434, 2012.
Corollary 5.2.
\url{https://arxiv.org/abs/1004.4389}

\bibitem{ST83}
Daniel D. Sleator and Robert E. Tarjan.
A data structure for dynamic trees.
\emph{Journal of Computer and System Sciences} 26(3), 362--391, 1983.
\url{https://doi.org/10.1016/0022-0000(83)90006-5}
\end{thebibliography}
\end{document}

